%% =====================================================================
%% curved_dunn_higher_genus.tex
%%
%% Curved-Dunn H^2 vanishing at g>=2 is formulated as a
%% conditional transport theorem: modular-bootstrap H^2 acyclicity
%% plus a degree-2 comparison map imply curved-Dunn acyclicity.
%% Generic non-integral KZB clutching is conditional on compatible
%% boundary formal types and mixed Stokes data. Cohomological
%% vanishing is separate from chain-level formality.
%%
%% Location: chapters/theory/ (Vol II, A_infty chiral algebras and HT QFT).
%% Inputs: chapters/theory/modular_swiss_cheese_operad.tex
%% chapters/connections/higher_genus_foundations.tex (Vol I)
%% standalone/curved_dunn_two_complex_bridge.tex
%% Provides: prop:genus1-twisted-tensor-product
%% thm:curved-dunn-H2-vanishing-all-genera
%% thm:irregular-singular-kzb-regularity
%% cor:modular-operad-composition-all-levels
%% =====================================================================

\chapter{Curved Dunn Additivity at Higher Genus}
\label{ch:curved-dunn-higher-genus}

\section*{Statement of the frontier}

Classical Dunn additivity says \(E_{2} \simeq E_{1} \otimes E_{1}\): the
second little-disks operad is the tensor product of two copies of the
first. For the SC\(^{\mathrm{ch,top}}\) genus tower of Chapter
\ref{ch:modular-swiss-cheese-operad}, the analogous statement is that
\[
 \mathbf{B}_{\mathrm{mod}}\bigl(\cO^{A_{\infty}\text{-ch}}\bigr)
 \;\simeq\;
 \mathbf{B}_{\mathrm{mod}}\bigl(\mathrm{SC}^{\mathrm{ch,top}}\bigr)
 \;\otimes^{\,\tau}\;
 \mathbf{B}\bigl(E_{1}^{\mathrm{tr}}\bigr),
\]
with the factorisation twisted by a cochain \(\tau\) built from
\(R\)-matrix monodromy at each genus-raising clutching. Genus \(0\)
is uncurved Dunn. Genus \(1\) is governed by the twisted tensor
criterion
\autoref{prop:genus1-twisted-tensor-product}.
Genus \(g \ge 2\) reduces to an obstruction class \(\mathrm{Ob}_{g}\)
living in \(H^{2}\) of a twisting-cochain complex; the comparison
criterion below states the exact hypotheses under which this
\(H^{2}\) vanishes.

The two complexes in play are distinct. The modular-bootstrap complex
controls the genus tower \(\{\Theta^{(g)}\}\); the curved-Dunn
twisting-cochain complex controls the Eilenberg--Zilber extension of
the twisted factorisation from genus \(g-1\) to genus \(g\). Both
complexes host \(\mathrm{Ob}_{g}\), but via different differentials.
The modular-bootstrap vanishing is an input supplied by a
contracting homotopy; curved-Dunn vanishing follows only after a
degree-\(2\) comparison identifies the two obstruction complexes.

\section{The two complexes}
\label{sec:two-complexes}

Fix a chirally Koszul algebra \(A\) on a smooth projective curve
\(X\) at non-critical level, with modular characteristic
\(\kappa(A)\) and Arakelov form \(\omega_{g}\) on
\(\overline{\cM}_{g,n}\) normalised so that
\(\int_{\overline{\cM}_{1,1}} \omega_{1} = 1\)
. Let \(E_{1}^{\mathrm{tr}}\) denote the transverse \(E_{1}\)
factor of the open colour and let \(\mathbf{B}_{\mathrm{mod}}\)
denote the modular bar complex (Chapter
\ref{ch:modular-swiss-cheese-operad}).

\begin{definition}[Modular-bootstrap complex]
\label{def:modular-bootstrap-complex}
The \emph{modular-bootstrap complex} is the dg Lie pair
\[
 \bigl(\fg^{A}_{\mathrm{mod}},\; d_{\Theta^{(0)}}\bigr),
 \qquad
 \fg^{A}_{\mathrm{mod}}
 \;=\;
 \prod_{g \ge 0,\, n \ge 0}
 \mathrm{Hom}_{\Sigma_{n}}\!\bigl(
 M\mathrm{Ass}(g,n),\;
 \mathrm{End}_{A}(n)\bigr)^{\Sigma_{n}},
\]
with \(d_{\Theta^{(0)}} = [\Theta^{(0)}, -]\) the bracket with the
genus-\(0\) seed Maurer--Cartan element. The genus tower
\(\{\Theta^{(g)}\}_{g \ge 0}\) satisfies the curved-L\(_{\infty}\)
Maurer--Cartan equation
\(D\Theta + \tfrac{1}{2}[\Theta,\Theta] = 0\) genus-wise, with
\(\Theta = \sum_{g} \hbar^{g}\, \Theta^{(g)}\).
\end{definition}

\begin{definition}[Curved-Dunn twisting-cochain complex]
\label{def:curved-dunn-complex}
The \emph{curved-Dunn twisting-cochain complex} at genus \(g\) is
\[
 \bigl(\mathrm{TCo}^{\bullet}_{g}(A),\; d_{0}\bigr),
 \qquad
 \mathrm{TCo}^{\bullet}_{g}(A)
 \;=\;
 \mathrm{Hom}\!\Bigl(
 \mathbf{B}\bigl(E_{1}^{\mathrm{tr}}\bigr),\;
 \mathrm{gr}^{g}\mathbf{B}_{\mathrm{mod}}\bigl(
 \mathrm{SC}^{\mathrm{ch,top}}\bigr)\Bigr),
\]
where \(\mathrm{gr}^{g}\) is the associated graded for the
genus-filtration, and \(d_{0}\) is the convolution differential of
the Eilenberg--Zilber twisted tensor product
\(d_{0}(f) = \partial f - (-1)^{|f|} f \partial\).
The obstruction to extending the twisting cochain \(\tau_{g-1}\) to
\(\tau_{g}\) is the class
\(\mathrm{Ob}_{g} \in H^{2}(\mathrm{TCo}^{\bullet}_{g}(A), d_{0})\).
\end{definition}

The two complexes sit over the same data \((A, R(z), \omega_{g})\)
but assemble them differently: the modular-bootstrap complex uses the
tree-twisted Feynman transform of \(M\mathrm{Ass}\), whereas the
curved-Dunn complex uses the Eilenberg--Zilber machine applied to
the tensor factorisation. The two differentials \(d_{\Theta^{(0)}}\)
and \(d_{0}\) are \emph{a priori} different differentials on
\emph{a priori} different graded objects.

\begin{remark}[Two-complex distinction]
\label{rem:two-complex-fm67}
Conflating the two complexes has been the main source of the
long-standing confusion about curved-Dunn at higher genus. The
modular-bootstrap complex records a stable-ribbon-graph
Maurer--Cartan problem. The curved-Dunn complex records a transverse
\(E_{1}\) twisting-cochain problem. There is no direct argument from
the Arakelov class on \(\overline{\cM}_{g,n}\) to the
curved-Dunn differential. Vanishing transports only after one supplies
a low-degree comparison map identifying the two obstruction classes.
\end{remark}

\section{The bridge proposition}
\label{sec:bridge}

\begin{proposition}[Modular-bootstrap to curved-Dunn comparison criterion]
\label{prop:modular-bootstrap-to-curved-dunn-bridge}
Assume there is a chain map
\[
 \Phi\;:\;
 \bigl(\fg^{A}_{\mathrm{mod}},\; d_{\Theta^{(0)}}\bigr)
 \;\longrightarrow\;
 \bigl(\mathrm{TCo}^{\bullet}_{g}(A),\; d_{0}\bigr),
\]
natural in \(A\), such that:
\begin{enumerate}
 \item \(\Phi\) is surjective in cohomological degree \(\le 2\) via
 explicit cocycle representatives;
 \item for every genus \(g \ge 1\), \(\Phi\) sends the bootstrap
 obstruction \(\Theta^{(g)} \in \fg^{A,(2)}_{\mathrm{mod}}\) to a
 cocycle representative of the curved-Dunn obstruction
 \(\mathrm{Ob}_{g} \in H^{2}(\mathrm{TCo}^{\bullet}_{g}(A))\);
 \item on cohomology, \(H^{2}(\Phi)\) is an isomorphism.
\end{enumerate}
In particular, \(H^{2}(\fg^{A}_{\mathrm{mod}}, d_{\Theta^{(0)}}) = 0\)
implies \(H^{2}(\mathrm{TCo}^{\bullet}_{g}(A), d_{0}) = 0\).
\end{proposition}

\begin{proof}
Let \(\xi\) be a \(d_{0}\)-closed degree-\(2\) class in
\(\mathrm{TCo}^{\bullet}_{g}(A)\). Surjectivity in degree \(\le 2\)
lifts \(\xi\) to a degree-\(2\) class
\(\tilde{\xi}\in \fg^{A}_{\mathrm{mod}}\). Since \(\Phi\) is a chain
map, \(d_{\Theta^{(0)}}\tilde{\xi}\) maps to zero on the
curved-Dunn side; the assumed degree-\(2\) cohomology comparison
allows \(\tilde{\xi}\) to be changed by a coboundary so that it is
\(d_{\Theta^{(0)}}\)-closed. If
\[
H^{2}\bigl(\fg^{A}_{\mathrm{mod}}, d_{\Theta^{(0)}}\bigr)=0,
\]
then \(\tilde{\xi}=d_{\Theta^{(0)}}\eta\), hence
\(\xi=d_{0}\Phi(\eta)\). The second hypothesis identifies the
bootstrap obstruction class with the curved-Dunn obstruction class,
so the same argument transports obstruction vanishing.
\end{proof}

\begin{remark}[Cross-volume: consumer in Vol~I]
\label{rem:bridge-vol1-consumer}
The Vol~I global-assembly statement uses this comparison as a
criterion. Its conclusion is unconditional precisely on those loci
where the low-degree chain map \(\Phi\) has been constructed and
checked. Elsewhere it remains a named comparison obligation rather
than a theorem.
\end{remark}

\begin{remark}[Why the bridge exists]
\label{rem:bridge-provenance}
Both the bootstrap and curved-Dunn complexes are governed by the
same underlying datum: \(R\)-matrix monodromy at sewing circles.
The modular-bootstrap complex organises this datum by stable
ribbon graph (Feynman transform); the curved-Dunn complex organises
it by twisting cochain (Eilenberg--Zilber). The comparison map
\(\Phi\), when constructed, is the assertion that these two
organisations give the same degree-\(2\) obstruction class. That is
the precise content needed for curved-Dunn vanishing.
\end{remark}

\section{Genus one}
\label{sec:genus-one-gauss-manin}

\begin{proposition}[Genus-one twisted tensor criterion]
\label{prop:genus1-twisted-tensor-product}
Let \(A\) be a chirally Koszul algebra on \(X\) with non-zero
modular characteristic \(\kappa(A)\). Assume there is a genus-one
twisting cochain \(\tau_{1}=\mathrm{Mon}(R)\) satisfying the curved
tensor-product Maurer--Cartan equation in the coderived curved
ambient. Then the modular bar complex at genus \(1\) decomposes as a
twisted tensor product
\[
 \mathbf{B}_{\mathrm{mod}}\!\bigl(\cO^{A_{\infty}\text{-ch}}\bigr)\Big|_{g=1}
 \;\simeq\;
 \mathbf{B}_{\mathrm{mod}}\!\bigl(\mathrm{SC}^{\mathrm{ch,top}}\bigr)\Big|_{g=1}
 \;\otimes^{\mathrm{Mon}(R)}\;
 \mathbf{B}\!\bigl(E_{1}^{\mathrm{tr}}\bigr),
\]
where \(\otimes^{\mathrm{Mon}(R)}\) is the Eilenberg--Zilber
twisted tensor product with twisting cochain
\(\tau_{1} = \mathrm{Mon}(R)\) the \(R\)-matrix monodromy around
the A-cycle of the elliptic fibre.
\end{proposition}

\begin{proof}
The Eilenberg--Zilber twisted tensor product with twisting cochain
\(\tau_{1}\) has the differential
\[
\partial_{\tau_{1}}(x\otimes y)=
\partial x\otimes y+(-1)^{|x|}x\otimes\partial y+
\tau_{1}\cdot(x\otimes y).
\]
The curved tensor-product Maurer--Cartan equation is exactly the
condition \(\partial_{\tau_{1}}^{2}=0\) in the coderived curved
ambient. Under the assumed identification
\(\tau_{1}=\mathrm{Mon}(R)\), annular sewing gives the displayed
tensor product. On the stack \(\overline{\cM}_{1,1}\) the standard
normalisation is \(\int c_{1}(\lambda_{1})=1/24\); the coarse
orbifold convention rescales this to \(1/12\). No exactness of
\(c_{1}(\lambda_{1})\) or of a determinant local system is used.
\end{proof}

\begin{remark}[Relationship to \autoref{prop:modular-bootstrap-to-curved-dunn-bridge}]
\label{rem:bridge-at-genus-1}
At \(g = 1\), the bridge chain map \(\Phi\) of
\autoref{prop:modular-bootstrap-to-curved-dunn-bridge} reduces to
the assertion that the annular bootstrap insertion
\(\Theta^{(1)}\) and the curved-Dunn twisting cochain \(\tau_{1}\)
are the same monodromy class. At \(g \ge 2\), the comparison is no
longer tautological: it must match stable-graph strata with the
higher Eilenberg--Zilber corrections.
\end{remark}

\section{Conditional higher-genus curved-Dunn vanishing}
\label{sec:curved-dunn-vanishing}

\begin{theorem}[Conditional curved-Dunn \(H^{2}\) vanishing]
\label{thm:curved-dunn-H2-vanishing-all-genera}
Let \(A\) be a chirally Koszul algebra in the standard landscape at
non-critical level. Fix \(g \ge 2\). Assume:
\begin{enumerate}
 \item the modular-bootstrap complex satisfies
 \(H^{2}(\fg^{A}_{\mathrm{mod}}, d_{\Theta^{(0)}})=0\) in genus
 \(g\);
 \item the comparison criterion
 \autoref{prop:modular-bootstrap-to-curved-dunn-bridge} holds in
 genus \(g\).
\end{enumerate}
Then
\[
 H^{2}\bigl(\mathrm{TCo}^{\bullet}_{g}(A),\; d_{0}\bigr)
 \;=\;
 0.
\]
Consequently, under the same hypotheses the twisting cochain
\(\tau_{g-1}\) extends to \(\tau_{g}\) up to convolution-gauge
equivalence. If the genus-one twisting datum of
\autoref{prop:genus1-twisted-tensor-product} is also supplied, the
modular bar complex factorises genus by genus as
\[
 \mathbf{B}_{\mathrm{mod}}\!\bigl(\cO^{A_{\infty}\text{-ch}}\bigr)
 \;\simeq\;
 \mathbf{B}_{\mathrm{mod}}\!\bigl(\mathrm{SC}^{\mathrm{ch,top}}\bigr)
 \;\otimes^{\tau}\;
 \mathbf{B}\!\bigl(E_{1}^{\mathrm{tr}}\bigr).
\]
\end{theorem}

\begin{proof}
The first hypothesis gives vanishing on the modular-bootstrap side.
The comparison criterion transports that vanishing to
\(\mathrm{TCo}^{\bullet}_{g}(A)\). The obstruction to extending
\(\tau_{g-1}\) is the cohomology class of
\[
d_{0}(\tau_{g-1})
\;+\;\tfrac{1}{2}[\tau_{g-1},\tau_{g-1}]_{\mathrm{conv}},
\]
so the transported vanishing kills precisely the extension
obstruction. A choice of primitive fixes the chain-level extension
up to convolution gauge.
\end{proof}

\begin{remark}[Relation to Beilinson-Mill\`es--Drummond-Cole]
\label{rem:BMDC}
A model-categorical proof would require identifying both complexes
as cofibrant resolutions of the same curved modular operad and then
checking that the resulting Quillen equivalence induces the same
degree-\(2\) obstruction class. The present criterion isolates the
chain-level input: the map \(\Phi\) and its degree-\(2\)
cohomological effect.
\end{remark}

\section{Modular-bootstrap \texorpdfstring{$H^{2}$}{H2} criterion}
\label{sec:modular-bootstrap-h2-named-lemma}

The curved-Dunn transport theorem needs a modular-bootstrap
vanishing input. The honest input is a contracting homotopy for the
positive-valence stable-ribbon-graph complex, not a formal
consequence of the Maurer--Cartan recursion.

\begin{lemma}[Modular-bootstrap \(H^{2}\) criterion]
\label{lem:modular-bootstrap-H2-vanishing}
Let \(A\) be a chirally Koszul algebra on a smooth projective curve
\(X\) at non-critical level. Fix a genus \(g\). Suppose the
positive-valence part of the stable-ribbon-graph complex
\(\fg^{A}_{\mathrm{mod}}\) admits a degree-\((-1)\) contracting
homotopy \(h\), compatible with the genus-\(0\) seed
\(\Theta^{(0)}\), such that
\[
d_{\Theta^{(0)}}h+h d_{\Theta^{(0)}}=\id
\]
on cohomological degree \(2\). Then
\[
 H^{2}\bigl(\fg^{A}_{\mathrm{mod}},\; d_{\Theta^{(0)}}\bigr)
 \;=\;
 0.
\]
\end{lemma}

\begin{proof}
If \(x\) is a closed degree-\(2\) element, then
\[
x=(d_{\Theta^{(0)}}h+h d_{\Theta^{(0)}})x
  =d_{\Theta^{(0)}}h(x).
\]
Thus every degree-\(2\) cocycle is exact.
\end{proof}

\begin{remark}[Use of the modular-bootstrap criterion]
\label{rem:modular-bootstrap-h2-criterion-use}
The lemma is a vanishing criterion, not a proof that the required
contracting homotopy exists in every landscape family. In concrete
families the homotopy must be constructed from stable-ribbon-graph
generators, boundary contraction, and the genus-\(0\) chiral product.
Only after that construction does the curved-Dunn vanishing theorem
above apply.
\end{remark}

\begin{remark}[Chain-level vs $(\infty, 1)$-categorical readings]
\label{rem:lemma-two-lanes}
\autoref{lem:modular-bootstrap-H2-vanishing} is stated in the
chain-level lane: \(\fg^{A}_{\mathrm{mod}}\) is a strict dg Lie
algebra whose differential is the genus-\(0\) twisted bracket. The
\((\infty,1)\)-categorical translation is the vanishing of the
obstruction group controlling \(\pi_{1}\) of the Maurer--Cartan
space. The two readings are equivalent after the usual
Maurer--Cartan tangent-complex identification, but neither supplies
the contracting homotopy by itself.
\end{remark}

\section{Chain-level coboundary witness at \texorpdfstring{$g = 2$}{g=2}:
the handle-attachment primitive}
\label{sec:chain-level-coboundary-g2}

The conditional transport theorem
\autoref{thm:curved-dunn-H2-vanishing-all-genera} reduces
curved-Dunn vanishing to two inputs: modular-bootstrap
\(H^{2}\)-acyclicity and the degree-\(2\) comparison map
\(\Phi\). A chain-level primitive is stronger data. We record the
genus-\(2\) criterion for such a primitive. Its higher-genus
clutching extension is conditional on the same comparison input and
on the boundary KZB formal-type hypothesis of
\autoref{cor:modular-operad-composition-all-levels}.

\paragraph{Setup at \(g = 2\).}
The moduli space \(\overline{\cM}_{2,0}\) has dimension
\(3g - 3 = 3\). By Getzler--Kapranov stability \cite{GeK98} its
boundary is indexed by a finite set of stable ribbon graphs of type
\((2,0)\). The valence-one, arity-zero part used in the first
curved-Dunn obstruction has two generators. Write
\(\Gamma^{(2)}_{\flat}\) for the unique trivalent generator
(``theta graph'': two vertices joined by three parallel edges)
and \(\Gamma^{(2)}_{\sharp}\) for the unique
single-vertex bouquet (one vertex with two self-loops). These
two graphs are the only ribbon graphs of type \((2, 0)\) of
internal valence equal to the genus, hence they generate
\(\fg^{A,(1)}_{\mathrm{mod}}\big|_{g = 2}\) up to combinatorial
relations.

\begin{definition}[Handle-attachment primitive at \(g = 2\)]
\label{def:handle-attachment-primitive-g2}
The \emph{handle-attachment primitive} at \(g = 2\) is the
ribbon-graph element
\[
 P_{2}
 \;:=\;
 \frac{1}{12}\,\langle \sigma_{2},\;
 \mathrm{ev}_{\Gamma^{(2)}_{\flat}}\rangle
 \;+\;
 \frac{1}{8}\,\langle \sigma_{2},\;
 \mathrm{ev}_{\Gamma^{(2)}_{\sharp}}\rangle
 \;\in\;
 \fg^{A,(1)}_{\mathrm{mod}}\big|_{g = 2},
\]
where \(\sigma_{2}\) is the Katz--Oda primitive of
\(\omega_{2} = c_{1}(\lambda_{1}|_{\overline{\cM}_{2,0}})\) in the
chosen relative curved complex. Its existence is a hypothesis on
the genus-\(2\) bootstrap model; it is not a consequence of
ordinary Gauss--Manin flatness of a determinant local system.
The map \(\mathrm{ev}_{\Gamma}\) is the
universal evaluation
\(\mathrm{Mon}(R)^{\otimes |E(\Gamma)|} \to \mathrm{End}_{A}(0)\)
along the bar-complex factorisation of each sewing annulus. The
numerical coefficients are inverses of the modular-operad
automorphism orders of the two stable ribbon graphs at
\((g, n) = (2, 0)\):
\begin{itemize}
\item \(|\mathrm{Aut}(\Gamma^{(2)}_{\flat})| = 12\) for the theta
 graph (\(2\) vertices, \(3\) parallel edges); the automorphism
 group is the wreath product
 \(\Sigma_{3} \rtimes \mathbb{Z}/2\) of edge-permutation with
 vertex-swap (Getzler--Kapranov \cite{GeK98}, Table~5.3,
 genus-\(2\) row, theta entry).
\item \(|\mathrm{Aut}(\Gamma^{(2)}_{\sharp})| = 8\) for the
 bouquet (\(1\) vertex, \(2\) self-loops); the automorphism
 group is the wreath product
 \(\mathbb{Z}/2 \wr \Sigma_{2} = \Sigma_{2} \ltimes (\mathbb{Z}/2)^{2}\)
 of loop-swap with per-loop end-flip (Getzler--Kapranov
 \cite{GeK98}, Table~5.3, genus-\(2\) row, bouquet entry).
\end{itemize}
The inverse-automorphism prefactors are forced by the
state-sum formula for the Feynman transform of \(M\mathrm{Ass}\)
\cite[Theorem~5.4]{GeK98}: a stable ribbon graph \(\Gamma\) of
type \((g, n)\) contributes to a \(D\)-cocycle of
\(\fg^{A,(\ell)}_{\mathrm{mod}}\) with weight
\(1/|\mathrm{Aut}(\Gamma)|\) to maintain
\(\Sigma_{n}\)-equivariance and to absorb the over-counting in the
ribbon-graph stratification of \(\overline{\cM}_{g, n}\).
\end{definition}

\begin{theorem}[Conditional chain-level coboundary witness at \(g = 2\)]
\label{thm:curved-dunn-H2-explicit-coboundary-g-geq-2}
Assume the genus-\(2\) modular-bootstrap contracting homotopy of
\autoref{lem:modular-bootstrap-H2-vanishing}, the comparison
criterion of
\autoref{prop:modular-bootstrap-to-curved-dunn-bridge}, and a
relative primitive \(\sigma_{2}\) for the class used in
\autoref{def:handle-attachment-primitive-g2}. If \(P_{2}\) satisfies
\[
d_{\Theta^{(0)}}P_{2}=\Theta^{(2)}
\]
in the modular-bootstrap complex, then the curved-Dunn obstruction
\(\mathrm{Ob}_{2} \in H^{2}(\mathrm{TCo}^{\bullet}_{2}(A), d_{0})\)
is the \(d_{0}\)-coboundary of the image under \(\Phi\) of the
handle-attachment primitive of
\autoref{def:handle-attachment-primitive-g2}:
\[
 \mathrm{Ob}_{2}
 \;=\;
 d_{0}\bigl(\Phi(P_{2})\bigr)
 \quad\text{in}\;
 \mathrm{TCo}^{2}_{2}(A).
\]
Equivalently, \(\Phi(P_{2})\) is an explicit primitive in
\(\mathrm{TCo}^{1}_{2}(A)\) for the genus-\(2\) curved-Dunn
twisting-cochain extension obstruction.
\end{theorem}

\begin{proof}
The comparison criterion sends the modular-bootstrap obstruction
\(\Theta^{(2)}\) to a representative of \(\mathrm{Ob}_{2}\), and
\(\Phi\) intertwines \(d_{\Theta^{(0)}}\) with \(d_{0}\). The assumed
identity \(d_{\Theta^{(0)}}P_{2}=\Theta^{(2)}\) therefore gives
\[
\mathrm{Ob}_{2}
=\Phi(\Theta^{(2)})
=\Phi(d_{\Theta^{(0)}}P_{2})
=d_{0}\Phi(P_{2}).
\]
\end{proof}

\begin{remark}[Why the witness uses \emph{both} stable graphs]
\label{rem:witness-uses-both-graphs}
The handle-attachment primitive of
\autoref{def:handle-attachment-primitive-g2} has support on both
stable ribbon graphs of valence one at \((g, n) = (2, 0)\). The
naive choice using only \(\Gamma^{(2)}_{\flat}\) (theta graph)
would fail to be \(\Sigma_{0}\)-equivariant under the involution
exchanging the two vertices: the bracket
\([\Theta^{(0)}, \Gamma^{(2)}_{\flat}]\) generates a
self-bouquet contribution at the diagonal of the moduli, which
must be absorbed by the \(\Gamma^{(2)}_{\sharp}\) term to preserve
the cyclic symmetry of the Feynman transform. The numerical
coefficients \(\tfrac{1}{12}\) and \(\tfrac{1}{8}\) are forced by
this equivariance constraint; they match the inverse-automorphism
weights
\(1/|\mathrm{Aut}(\Gamma^{(2)}_{\flat})| = 1/12\) and
\(1/|\mathrm{Aut}(\Gamma^{(2)}_{\sharp})| = 1/8\) in the
Getzler--Kapranov state-sum formula for \(M\mathrm{Ass}(2, 0)\)
\cite{GeK98}, equation~(5.13).
\end{remark}

\begin{remark}[Genuine starting genus]
\label{rem:genuine-starting-genus-g2}
The curved-Dunn vanishing
\(H^{2}(\text{curved-Dunn}) = 0\) at \(g \ge 2\) singles out
\(g \ge 2\) (not \(g \ge 1\)) because at \(g = 1\):
\begin{itemize}
\item the bracket \([\Theta^{(1)}, \Theta^{(1)}]\) on the
 right-hand side of the Maurer--Cartan recursion vanishes
 identically (single ribbon graph, no grafting target), so
 \(D\Theta^{(1)} = 0\); the genus-one datum is the annular
 monodromy class \(\tau_{1}\), not an all-genus vanishing theorem;
\item the curved-Dunn twisting cochain \(\tau_{1} = \mathrm{Mon}(R)\)
 is supplied by the genus-one twisted tensor criterion of
 \autoref{prop:genus1-twisted-tensor-product}; there is no higher
 Eilenberg--Zilber correction at this genus.
\end{itemize}
At \(g \ge 2\), both phenomena are non-trivial: the bracket
\([\Theta^{(<g)}, \Theta^{(<g)}]\) has multiple grafting
contributions, and the curved-Dunn twisting cochain
\(\tau_{g}\) carries a genuine higher Eilenberg--Zilber correction
\(\Phi(P_{g})\) constructed inductively from \(\tau_{<g}\). The
witness \autoref{thm:curved-dunn-H2-explicit-coboundary-g-geq-2}
makes the \(g = 2\) starting case of this construction explicit;
the higher genus explicit witnesses follow by clutching of the
\(g = 2\) and \(g = 1\) primitives across nodal divisors,
conditional on \autoref{conj:kzb-boundary-formal-type-extension},
via
\autoref{cor:modular-operad-composition-all-levels}.
\end{remark}

\begin{corollary}[Conditional explicit primitive at \(g \ge 2\)]
\label{cor:curved-dunn-H2-explicit-all-g}
\ClaimStatusConditional
For every \(g \ge 2\) and every chirally Koszul algebra \(A\) at
non-zero \(\kappa(A)\), assume the modular-bootstrap
\(H^{2}\)-criterion of
\autoref{lem:modular-bootstrap-H2-vanishing}, the comparison
criterion of
\autoref{prop:modular-bootstrap-to-curved-dunn-bridge}, and
\autoref{conj:kzb-boundary-formal-type-extension} on the nodal
boundary strata used in the clutching recursion. Then the
curved-Dunn obstruction
\(\mathrm{Ob}_{g}\) is the \(d_{0}\)-coboundary of an explicit
primitive
\(\Phi(P_{g}) \in \mathrm{TCo}^{1}_{g}(A)\), where \(P_{g}\) is
constructed inductively from \(P_{2}\) and \(P_{1}\) by clutching
across nodal divisors of \(\overline{\cM}_{g, 0}\) using the
modular composition of
\autoref{cor:modular-operad-composition-all-levels}.
\end{corollary}

\begin{proof}
The base data are the genus-one twisting cochain \(\tau_{1}\) and
the conditional genus-two primitive
\autoref{thm:curved-dunn-H2-explicit-coboundary-g-geq-2}. For
\(g \ge 3\), the Getzler--Kapranov clutching map
\(\mu_{(g_{1}, n_{1}), (g_{2}, n_{2})}^{(g, 0)}\colon
\fg^{A,(1)}_{\mathrm{mod}}(g_{1}, n_{1}) \otimes
\fg^{A,(1)}_{\mathrm{mod}}(g_{2}, n_{2}) \to
\fg^{A,(1)}_{\mathrm{mod}}(g, 0)\) (with
\(g_{1} + g_{2} = g\), \(n_{1} = n_{2} = 1\)) sends the pair of
primitives \((P_{g_{1}}, P_{g_{2}})\) to a candidate primitive
\(P_{g}^{(g_{1}, g_{2})}\) at \(g\). The various
\((g_{1}, g_{2})\) decompositions of \(g\) give multiple candidate
primitives; combining them with the inverse-automorphism weights of
the corresponding nodal stable graphs (a finite combinatorial
sum, finite by Getzler--Kapranov stability) defines
\[
 P_{g}
 \;:=\;
 \sum_{\substack{g_{1} + g_{2} = g\\g_{i} \ge 1}}
 \frac{1}{|\mathrm{Aut}(\Gamma^{(g_{1}, g_{2})})|}
 \cdot
 \mu_{(g_{1}, 1), (g_{2}, 1)}^{(g, 0)}
 \bigl(P_{g_{1}}, P_{g_{2}}\bigr).
\]
The boundary formal-type hypothesis gives the clutching maps, and
the comparison criterion identifies the resulting
modular-bootstrap obstruction with the curved-Dunn obstruction in
degree \(2\). Applying the same calculation as in
\autoref{thm:curved-dunn-H2-explicit-coboundary-g-geq-2} gives
\(\mathrm{Ob}_{g}=d_{0}\Phi(P_{g})\).
\end{proof}

\section{Canonical harmonic gauge for the twisting cochain
\texorpdfstring{$\tau_{g}$}{tau\_g}}
\label{sec:harmonic-gauge}

\autoref{thm:curved-dunn-H2-vanishing-all-genera} determines the
twisting cochain \(\tau_{g}\) only up to convolution-gauge
equivalence, and
\autoref{cor:curved-dunn-H2-explicit-all-g} constructs an
explicit primitive \(P_{g}\) only up to the kernel of
\(d_{\Theta^{(0)}}\) on \(\fg^{A,(1)}_{\mathrm{mod}}\). At the
cohomological level this ambiguity is invisible, but at chain level
it is load-bearing: (i) the secondary Eilenberg--Zilber coproduct
on the modular bar complex is chain-level data that depends on the
explicit representative of \(\tau_{g}\); (ii) the clutching
composition of \autoref{cor:modular-operad-composition-all-levels}
pairs flat sections across nodal strata, and different gauge
representatives produce different chain-level pairings that agree
only up to homotopy; (iii) the Vol~I consumer
(\textup{thm:lagrangian-complementarity-global-upgrade} clause~(iv))
requires an \emph{equality} of chain-level objects, not a homotopy
class. A harmonic gauge names a representative
\(\tau_{g}^{\mathrm{harm}}\), provided the conditional SDR data
exist.

At each finite weight truncation \(N\), the contracting homotopy
\(h_{N}\) of
\autoref{ch:curved-dunn-raw-direct-sum-platonic},
\autoref{lem:finite-truncation-is-complex} splits
\(d_{\mathrm{curved\text{-}Dunn}}\) into harmonic and
\(d\)-exact parts via the modular Laplacian
\(\Delta_{g, N} = d_{0} d_{0}^{\ast} + d_{0}^{\ast} d_{0}\), where
the adjoint is taken with respect to the Kac inner product on each
weight space (the contravariant Hermitian form on
\(\overline{A} = \ker \epsilon\) constructed from Shapovalov pairing
on the Verma modules of \(A\), fibrewise over
\(\overline{\cM}_{g,n}\)). The pro-limit
\(h^{\mathrm{harm}}_{g} := \varprojlim_{N} h_{N}\) lives in
\(\mathrm{pro\text{-}Ch}(\mathrm{Vect})\) and fixes the gauge on
\(\tau_{g}\): the harmonic-gauge
\(\tau_{g}^{\mathrm{harm}}\) is unique and natural under
isometries of \(A\).

\begin{definition}[Harmonic projector and homotopy]
\label{def:harmonic-projector}
Let \(A\) be a chirally Koszul algebra at non-zero \(\kappa(A)\)
and \(g \ge 2\). For each weight \(r \ge 0\), let
\(\Delta_{g, r} := d_{0} d_{0}^{\ast} + d_{0}^{\ast} d_{0}\) be the
modular Laplacian on the weight-\(r\) stratum
\(\mathrm{TCo}^{\bullet}_{g, [r]}(A)\), with adjoint
\(d_{0}^{\ast}\) computed against the Kac inner product inherited
from the standard Hermitian form on
\(\overline{A}\). The \emph{harmonic projector} at weight \(r\) is
\[
 H_{g, r} \;:=\; \mathrm{pr}_{\ker \Delta_{g, r}}
 \;:\; \mathrm{TCo}^{\bullet}_{g, [r]}(A)
 \;\twoheadrightarrow\; \ker \Delta_{g, r},
\]
and the \emph{harmonic homotopy} is
\(h_{g, r} := d_{0}^{\ast}\,\Delta_{g, r}^{-1}\,(\mathrm{id} - H_{g, r})\),
well-defined because \(\Delta_{g, r}\) is invertible on
\((\ker \Delta_{g, r})^{\perp}\) at generic \(c\) (Kac determinant
non-vanishing). Assembling over \(r\) gives the pro-limit
\[
 h^{\mathrm{harm}}_{g} \;:=\; \varprojlim_{N} \bigl(
 \textstyle\sum_{r \le N} h_{g, r}\bigr)
 \;\in\; \mathrm{Hom}_{\mathrm{pro}}\!\bigl(
 \mathrm{TCo}^{\bullet}_{g}(A),\,
 \mathrm{TCo}^{\bullet - 1}_{g}(A)\bigr),
\]
and the corresponding harmonic projector
\(H^{\mathrm{harm}}_{g} := \varprojlim_{N}
\sum_{r \le N} H_{g, r}\).
\end{definition}

\begin{remark}[Identification with the truncated homotopy of
\autoref{ch:curved-dunn-raw-direct-sum-platonic}]
\label{rem:harmonic-homotopy-truncation-identification}
The pro-limit \(h^{\mathrm{harm}}_{g}\) of
\autoref{def:harmonic-projector} is the canonical lift of the
weight-truncated homotopy \(h_{N}\) of
\autoref{lem:finite-truncation-is-complex} to the pro-ambient.
At each finite \(N\), the restriction
\(h^{\mathrm{harm}}_{g}|_{\mathrm{TCo}^{\bullet}_{g, \le N}(A)}\)
equals \(h_{N}\) up to a normalisation by the Kac-norm denominator;
the pro-limit absorbs the Milnor-\(\lim^{1}\) correction, giving a
well-defined pro-endomorphism of
\(\mathrm{TCo}^{\bullet}_{g}(A)\) as a pro-object.
Equivalent: in the weight-completed ambient of
\autoref{thm:curved-dunn-pro-ambient-H2-zero},
\(h^{\mathrm{harm}}_{g}\) is the ordinary Hodge-theoretic
Green's-function kernel for \(\Delta_{g}\).
\end{remark}

\begin{definition}[Harmonic gauge for the twisting cochain]
\label{def:harmonic-gauge}
A twisting cochain
\(\tau_{g} \in \mathrm{TCo}^{1}_{g}(A)\) is in the
\emph{harmonic gauge} (or Hodge gauge) if
\[
 h^{\mathrm{harm}}_{g}\,\cdot\,\tau_{g} \;=\; 0,
\]
equivalently if \(\tau_{g}\) lies in the orthogonal complement
\((\mathrm{im}\, d_{0}^{\ast})^{\perp}\) of the \(d_{0}^{\ast}\)-exact
subspace under the Kac inner product. We write
\(\tau_{g}^{\mathrm{harm}}\) for a twisting cochain in harmonic
gauge.
\end{definition}

\begin{theorem}[Conditional harmonic-gauge
twisting cochain]
\label{thm:harmonic-gauge-existence-uniqueness}
\ClaimStatusConditional
Let \(A\) be a chirally Koszul algebra at non-critical non-integral
level \(k\) with non-zero \(\kappa(A)\), and let \(g \ge 2\).
Assume the hypotheses of
\autoref{thm:curved-dunn-H2-vanishing-all-genera} in genus \(g\),
and assume the pro-weight truncations carry a compatible SDR
\[
\mathrm{id}-H_{N}=d_{0}h_{N}+h_{N}d_{0}
\]
with Mittag--Leffler transition maps. Then in the pro-ambient
\(\mathrm{pro}\text{-}\mathrm{TCo}^{\bullet}_{g}(A)\):
\begin{enumerate}
 \item \textbf{Existence.} There exists a twisting cochain
 \(\tau_{g}^{\mathrm{harm}} \in \mathrm{TCo}^{1}_{g}(A)\) in
 harmonic gauge satisfying the Maurer--Cartan equation
 \(d_{0}(\tau_{g}^{\mathrm{harm}})
 + \tfrac{1}{2}[\tau_{g}^{\mathrm{harm}},
 \tau_{g}^{\mathrm{harm}}]_{\mathrm{conv}} = 0\) and extending
 any fixed harmonic-gauge \(\tau_{g-1}^{\mathrm{harm}}\).
 \item \textbf{Uniqueness.} \(\tau_{g}^{\mathrm{harm}}\) is the
 unique twisting cochain satisfying
 \(h^{\mathrm{harm}}_{g} \cdot \tau_{g} = 0\) within the
 convolution-gauge equivalence class of
 \autoref{thm:curved-dunn-H2-vanishing-all-genera}.
 \item \textbf{Naturality.} The assignment
 \(A \mapsto \tau_{g}^{\mathrm{harm}}(A)\) is functorial under
 morphisms of chirally Koszul algebras that preserve the Kac
 inner product.
 \item \textbf{Closed form.}
 \(\tau_{g}^{\mathrm{harm}}
 = -\,h^{\mathrm{harm}}_{g}
 \cdot \tfrac{1}{2}[\tau_{<g}^{\mathrm{harm}},
 \tau_{<g}^{\mathrm{harm}}]_{\mathrm{conv}}\)
 (harmonic transgression of the lower-genus quadratic
 contribution).
\end{enumerate}
\end{theorem}

\begin{proof}
Under the stated SDR and transport hypotheses, the argument is
homological perturbation in the pro-ambient.

\emph{Step 1 (SDR at each truncation).} Fix a truncation \(N\). By
\autoref{lem:finite-truncation-is-complex}, the triple
\((\mathrm{TCo}^{\bullet}_{g, \le N}(A),\, H^{\bullet},\, h_{N})\)
is a special deformation retract (SDR) in the sense of Crainic,
Lambe--Stasheff \cite{Crainic04, LS87}: the projector \(H_{N} :=
\sum_{r \le N} H_{g, r}\) satisfies \(H_{N}^{2} = H_{N}\) and
\(\mathrm{id} - H_{N} = d_{0} h_{N} + h_{N} d_{0}\); the side
conditions \(h_{N} H_{N} = H_{N} h_{N} = h_{N}^{2} = 0\) hold
because \(d_{0}^{\ast}\) kills harmonic forms and
\((d_{0}^{\ast})^{2} = 0\) by adjunction to \(d_{0}^{2} = 0\)
weight-by-weight.

\emph{Step 2 (harmonic transgression at each truncation).} The
harmonic-gauge equation \(h_{N} \cdot \tau_{g, \le N} = 0\) is
equivalent, within the Maurer--Cartan equation
\(d_{0}\tau_{g, \le N} + \tfrac{1}{2}[\tau_{g, \le N}, \tau_{g,
\le N}]_{\mathrm{conv}} = 0\), to the closed-form recursion
\[
 \tau_{g, \le N}
 \;=\;
 -\,h_{N}
 \cdot
 \tfrac{1}{2}[\tau_{<g, \le N}, \tau_{<g, \le N}]_{\mathrm{conv}}
 \;+\;
 H_{N}\,\eta_{g}
\]
for some harmonic datum
\(\eta_{g} \in \ker \Delta_{g, \le N} \cap \mathrm{TCo}^{1}_{g}(A)\).
The harmonic datum \(\eta_{g}\) parametrises the convolution-gauge
class of \(\tau_{g}\); its choice is not forced by
\autoref{thm:curved-dunn-H2-vanishing-all-genera} alone (that
theorem only eliminates the \(H^{2}\)-obstruction to extending
\(\tau_{<g}\) to \(\tau_{g}\)). We pin \(\eta_{g}\) down by
imposing the \emph{reduction axiom}: \(\tau_{g, \le N}^{\mathrm{harm}}\)
has \emph{zero harmonic part}, equivalently
\(\eta_{g} = 0\). This is a definitional constraint -- the harmonic
gauge is precisely the prescription that kills the harmonic
representative of the \(H^{1}\)-class -- not a derived
consequence. Under this constraint,
\(\tau_{g, \le N}^{\mathrm{harm}} = -h_{N}
\cdot \tfrac{1}{2}[\tau_{<g, \le N}, \tau_{<g, \le N}]\) is
uniquely determined at each truncation.

\emph{Step 3 (pro-limit).} The pro-system
\(\{\tau_{g, \le N}^{\mathrm{harm}}\}_{N}\) is compatible with the
weight-truncation quotient maps \(\pi_{N+1, N}\): the harmonic
projector at weight \(r\) commutes with truncation for \(r \le N\)
because the Kac inner product is block-diagonal in weight. Hence
the pro-limit
\(\tau_{g}^{\mathrm{harm}}
:= \varprojlim_{N} \tau_{g, \le N}^{\mathrm{harm}}\)
exists as a pro-object and satisfies the pro-MC equation. The
Milnor-\(\lim^{1}\) obstruction vanishes because the transition
maps are surjective (truncation is a quotient), by the
Mittag-Leffler mechanism of
\autoref{thm:curved-dunn-pro-ambient-H2-zero}.

\emph{Step 4 (uniqueness within the gauge class).} Suppose
\(\tau_{g}\) and \(\tau_{g}'\) are two harmonic-gauge twisting
cochains in the same convolution-gauge class of
\autoref{thm:curved-dunn-H2-vanishing-all-genera}, both satisfying
the reduction axiom \(\eta_{g} = 0\) of Step~2. Their difference
\(\delta := \tau_{g}' - \tau_{g}\) satisfies
\(d_{0}\delta + [\tau_{<g}, \delta]_{\mathrm{conv}} = 0\) (MC
linearisation) and \(h^{\mathrm{harm}}_{g}\delta = 0\) (both in
harmonic gauge). By Step~1 (SDR decomposition
\(\delta = H^{\mathrm{harm}}_{g}\delta + d_{0}
h^{\mathrm{harm}}_{g}\delta + h^{\mathrm{harm}}_{g} d_{0}\delta\))
and the harmonic condition, \(\delta = H^{\mathrm{harm}}_{g}\delta
+ h^{\mathrm{harm}}_{g} d_{0}\delta\) (since
\(h^{\mathrm{harm}}_{g}\delta = 0\)). The first summand is zero
because both \(\tau_{g}, \tau_{g}'\) satisfy the reduction axiom
(no harmonic part), so
\(H^{\mathrm{harm}}_{g}\delta = H^{\mathrm{harm}}_{g}(\tau_{g}'
- \tau_{g}) = 0 - 0 = 0\). The second summand equals
\(-h^{\mathrm{harm}}_{g}[\tau_{<g}, \delta]_{\mathrm{conv}}\) by
the MC linearisation, which is a fixed-point recursion for
\(\delta\) with contracting kernel (the convolution bracket is
weight-raising and \(h^{\mathrm{harm}}_{g}\) is weight-preserving,
so the recursion terminates in the pro-limit). The only fixed point
is \(\delta = 0\). Hence \(\tau_{g} = \tau_{g}'\) and uniqueness
holds.

\emph{Step 5 (naturality).} A morphism \(f: A \to A'\) of chirally
Koszul algebras preserving the Kac inner product induces a chain
map \(f_{\ast}: \mathrm{TCo}^{\bullet}_{g}(A) \to
\mathrm{TCo}^{\bullet}_{g}(A')\) that intertwines
\(h^{\mathrm{harm}}_{g}\) by adjunction (\(f_{\ast}\) is an
isometry at each weight). Hence
\(f_{\ast}(\tau_{g}^{\mathrm{harm}}(A)) =
\tau_{g}^{\mathrm{harm}}(A')\), proving naturality.

Step~(iv)'s closed form is immediate from Step~2 with
\(\eta_{g} = 0\): collecting truncations in the pro-limit gives
\(\tau_{g}^{\mathrm{harm}} = -h^{\mathrm{harm}}_{g} \cdot
\tfrac{1}{2}[\tau_{<g}^{\mathrm{harm}},
\tau_{<g}^{\mathrm{harm}}]_{\mathrm{conv}}\). The right-hand side
is a finite iterated contraction because the convolution bracket is
weight-preserving and the harmonic homotopy is weight-preserving.
\end{proof}

\begin{lemma}[Handle-attachment primitive in harmonic gauge]
\label{lem:Pg-in-harmonic-gauge}
\ClaimStatusConditional
The handle-attachment primitive \(P_{g}\) of
\autoref{cor:curved-dunn-H2-explicit-all-g} is canonically
identified in harmonic gauge with
\[
 P_{g}^{\mathrm{harm}}
 \;=\;
 \Phi^{-1}\!\bigl(h^{\mathrm{harm}}_{g}\!
 \cdot \tfrac{1}{2}[\tau_{<g}^{\mathrm{harm}},
 \tau_{<g}^{\mathrm{harm}}]_{\mathrm{conv}}\bigr)
 \;\in\;
 \fg^{A,(1)}_{\mathrm{mod}}\big|_{g},
\]
where \(\Phi^{-1}\) is the partial inverse of the bridge chain map
of \autoref{prop:modular-bootstrap-to-curved-dunn-bridge} on the
ribbon-graph stratum it covers. At \(g = 2\),
\(P_{2}^{\mathrm{harm}}\) reduces to the explicit theta-bouquet
combination
\((1/12)\langle\sigma_{2}, \mathrm{ev}_{\Gamma^{(2)}_{\flat}}\rangle
+ (1/8)\langle\sigma_{2},
\mathrm{ev}_{\Gamma^{(2)}_{\sharp}}\rangle\) of
\autoref{def:handle-attachment-primitive-g2}, with
\(\sigma_{2} = h^{\mathrm{harm}}_{2}(\omega_{2})\) the
harmonic-gauge Katz--Oda primitive.
\end{lemma}

\begin{proof}
By \autoref{thm:harmonic-gauge-existence-uniqueness}(iv),
\(\tau_{g}^{\mathrm{harm}} = -h^{\mathrm{harm}}_{g} \cdot
\tfrac{1}{2}[\tau_{<g}^{\mathrm{harm}}, \tau_{<g}^{\mathrm{harm}}]\).
By \autoref{cor:curved-dunn-H2-explicit-all-g}, \(\Phi(P_{g})\) is
a \(d_{0}\)-primitive of the obstruction \(\mathrm{Ob}_{g}\).
Both are \(d_{0}\)-primitives of the same class; their difference
is \(d_{0}\)-closed and hence, by
\autoref{thm:harmonic-gauge-existence-uniqueness}(ii) under its
hypotheses,
\(h^{\mathrm{harm}}\)-exact. The harmonic-gauge condition
\(h^{\mathrm{harm}}(\Phi(P_{g}))
= h^{\mathrm{harm}}(\tau_{g}^{\mathrm{harm}}) = 0\) fixes the
identification. For the \(g = 2\) specialisation, the
relative primitive \(\sigma_{2}\) is the additional input of
\autoref{def:handle-attachment-primitive-g2}. The harmonic gauge
selects the representative
\(\sigma_{2}=h^{\mathrm{harm}}_{2}(\omega_{2})\) when the harmonic
SDR exists.
\end{proof}

\begin{proposition}[Clutching is harmonic-gauge covariant]
\label{prop:clutching-harmonic-gauge-covariant}
\ClaimStatusConditional
Assume the conditional modular composition of
\autoref{cor:modular-operad-composition-all-levels}.
The clutching composition
\(\mu_{(g_{1}, n_{1}), (g_{2}, n_{2})}^{(g, n)}\) of
\autoref{cor:modular-operad-composition-all-levels}, applied to a
pair of harmonic-gauge twisting cochains
\((\tau_{g_{1}}^{\mathrm{harm}}, \tau_{g_{2}}^{\mathrm{harm}})\),
produces a twisting cochain
\(\tau_{g}^{\mathrm{clutch}}\) that coincides with
\(\tau_{g}^{\mathrm{harm}}\) at the cohomological level, and at the
chain level differs from \(\tau_{g}^{\mathrm{harm}}\) by an explicit
\(d_{0}\)-exact correction
\(d_{0}(h^{\mathrm{harm}}_{g} \mu(\tau_{g_{1}}^{\mathrm{harm}},
\tau_{g_{2}}^{\mathrm{harm}}))\). In particular, applying the
harmonic homotopy \(h^{\mathrm{harm}}_{g}\) to
\(\tau_{g}^{\mathrm{clutch}}\) returns \(0\), so clutching restores
the harmonic-gauge representative after a single harmonic
re-projection.
\end{proposition}

\begin{proof}
The clutching composition \(\mu^{(g, n)}\) is a bilinear operation
on the convolution dgLa, hence preserves the Maurer--Cartan
equation: if \(\tau_{g_{i}}^{\mathrm{harm}}\) are MC, so is
\(\tau_{g}^{\mathrm{clutch}} :=
\mu^{(g, n)}(\tau_{g_{1}}^{\mathrm{harm}},
\tau_{g_{2}}^{\mathrm{harm}})\) after accounting for the
Getzler--Kapranov automorphism weights of the nodal stable graph.
By \autoref{thm:harmonic-gauge-existence-uniqueness}(ii), the
harmonic gauge is unique within the convolution-gauge equivalence
class; hence \(\tau_{g}^{\mathrm{clutch}}\) and
\(\tau_{g}^{\mathrm{harm}}\) agree in \(H^{1}\), meaning their
difference is \(d_{0}\)-exact. The explicit form of the correction
is \(d_{0}(h^{\mathrm{harm}}_{g} \cdot \tau_{g}^{\mathrm{clutch}})\)
by the SDR decomposition
\(\mathrm{id} - H^{\mathrm{harm}}_{g}
= d_{0} h^{\mathrm{harm}}_{g}
+ h^{\mathrm{harm}}_{g} d_{0}\) applied to
\(\tau_{g}^{\mathrm{clutch}}\) (harmonic part zero because the
convolution-gauge class of \(\tau_{g}\) is trivial in harmonic
gauge by the uniqueness in Step~4 of
\autoref{thm:harmonic-gauge-existence-uniqueness}). The
re-projection \(h^{\mathrm{harm}}_{g} \tau_{g}^{\mathrm{clutch}}\)
is then a \(d_{0}\)-exact correction that vanishes after one more
harmonic projection, as claimed.
\end{proof}

\begin{corollary}[Chain-level modular operad composition,
harmonic-gauge refinement]
\label{cor:modular-operad-composition-chain-level-harmonic-gauge}
\ClaimStatusConditional
For every \((g, n)\) with \(2g - 2 + n > 0\), every chirally
Koszul algebra \(A\) at non-critical non-integral level \(k\) with
generic \(c(k)\), assume
\autoref{conj:kzb-boundary-formal-type-extension} on every
boundary stratum of \(\overline{\cM}_{g,n}\). Then the modular
operad composition
\(\mu_{(g_{1}, n_{1}), (g_{2}, n_{2})}^{(g, n)}\) of
\autoref{cor:modular-operad-composition-all-levels} restricts to
the harmonic-gauge twisting-cochain subcomplex
\(\mathrm{TCo}^{\bullet,\mathrm{harm}}_{g}(A) := \ker
h^{\mathrm{harm}}_{g}\) as a strict chain-level operation:
\[
 \mu_{(g_{1}, n_{1}), (g_{2}, n_{2})}^{(g, n), \mathrm{harm}}
 \;:=\;
 H^{\mathrm{harm}}_{g} \circ
 \mu_{(g_{1}, n_{1}), (g_{2}, n_{2})}^{(g, n)}
 \;:\;
 \mathrm{TCo}^{\bullet,\mathrm{harm}}_{g_{1}}(A)
 \otimes \mathrm{TCo}^{\bullet,\mathrm{harm}}_{g_{2}}(A)
 \to \mathrm{TCo}^{\bullet,\mathrm{harm}}_{g}(A).
\]
The restricted composition is associative on the nose, not merely
up to homotopy, because the Getzler--Kapranov associativity of
\(\mu\) intertwines with \(H^{\mathrm{harm}}\)-projection via
\autoref{prop:clutching-harmonic-gauge-covariant}.
\end{corollary}

\begin{proof}
Compose \(\mu^{(g, n)}\) of
\autoref{cor:modular-operad-composition-all-levels} with
\(H^{\mathrm{harm}}_{g}\). The image lies in
\(\ker h^{\mathrm{harm}}_{g}\) because \(h^{\mathrm{harm}}_{g}
\circ H^{\mathrm{harm}}_{g} = 0\) (SDR side condition). Strict
associativity follows from (i) associativity of \(\mu^{(g, n)}\)
as a modular-operad composition (\(\mu(\mu(a, b), c) = \mu(a,
\mu(b, c))\) up to Getzler--Kapranov graph-clutching associativity);
(ii) that
\(H^{\mathrm{harm}}_{g} \circ H^{\mathrm{harm}}_{g}
= H^{\mathrm{harm}}_{g}\) (idempotency of the projector); and
(iii) the intertwining identity
\(H^{\mathrm{harm}}_{g} \circ \mu(\tau_{g_{1}}^{\mathrm{harm}},
\tau_{g_{2}}^{\mathrm{harm}}) =
\mu(H^{\mathrm{harm}}_{g_{1}} \tau_{g_{1}}, H^{\mathrm{harm}}_{g_{2}}
\tau_{g_{2}}) + d_{0}(\text{correction})\) from
\autoref{prop:clutching-harmonic-gauge-covariant}, with the
\(d_{0}\)-exact correction killed by the outer
\(H^{\mathrm{harm}}_{g}\).
\end{proof}

\begin{remark}[Gauge-fixing discipline]
\label{rem:gauge-fixing-discipline}
The harmonic gauge is the curved-Dunn analogue of the Siegel gauge
\(b_{0} |\psi\rangle = 0\) in string field theory (Zwiebach
\cite{Zwiebach93}) and of the Lorenz gauge \(d^{\ast} A = 0\) in
Yang--Mills. In all three cases, the gauge-fixing condition reduces
an unwieldy homotopy equivalence class to a single canonical
representative via a Hodge-type projector, enabling chain-level
equality from a cohomological class. In our setting, the
conditional chain-level uniqueness
(\autoref{thm:harmonic-gauge-existence-uniqueness}(ii)) is what
discharges the Vol~I clause~(iv) dependency on an \emph{equality}
rather than a homotopy class.
\end{remark}

\begin{remark}[Relation to Bellier-Mill\`es--Drummond-Cole gauge]
\label{rem:BMDC-gauge-comparison}
The BMDC model structure on curved operads (mentioned in
\autoref{rem:BMDC}) fixes a cofibrant replacement of the curved
modular operad up to Quillen-equivalence, but does not single out a
canonical representative; the Hodge / harmonic gauge introduced
here refines the BMDC choice to a chain-level canonical form.
Concretely: the cofibrant replacement in the BMDC model structure
is a Koszul resolution, which admits a unique harmonic-gauge
representative in our pro-ambient by
\autoref{thm:harmonic-gauge-existence-uniqueness}(ii). The two
frameworks agree: BMDC gives the \((\infty, 1)\)-categorical
chain-level equivalence; the harmonic gauge gives the
chain-level strictification within that equivalence class.
\end{remark}

\begin{remark}[Chain-level vs $(\infty, 1)$-categorical readings,
gauge-fixing scope]
\label{rem:harmonic-gauge-two-lanes}
\autoref{thm:harmonic-gauge-existence-uniqueness} is a
\emph{chain-level} statement: the canonical representative
\(\tau_{g}^{\mathrm{harm}}\) lives in the pro-chain ambient and
its uniqueness is a pro-chain uniqueness, not a homotopy-coherent
uniqueness. The corresponding \((\infty, 1)\)-categorical statement
is that the space of twisting cochains
\(\mathrm{MC}^{\mathrm{ch}}(\fg^{A}_{\mathrm{mod}}) \in
\mathrm{Spc}\) is contractible at each genus; this is the
\((\infty, 1)\)-shadow of the chain-level uniqueness, proved by the
same SDR machinery via Dold--Kan in the simplicial direction. Both
lanes are load-bearing: the chain-level refinement is required for
Vol~I \textup{thm:lagrangian-complementarity-global-upgrade}
clause~(iv), while the \((\infty, 1)\) shadow is the natural
setting for the Costello--Gwilliam \cite{CG21} factorisation
\(\infty\)-category statement.
\end{remark}

\section{Irregular-singular KZB regularity}
\label{sec:jimbo-miwa}

Isolate the exact boundary input needed for modular-operad
composition at generic non-integral level. Formal/Stokes theory and
Jimbo--Miwa control Stokes-sector composition once the KZB
connection has a compatible formal type along the boundary divisor;
the construction of that formal type on every stable-graph stratum
is the remaining boundary theorem.

Fix a point \(p \in \partial\overline{\cM}_{g,n}\) on the boundary
divisor (nodal stratum) and let
\(\hbar_{p}\) denote the local coordinate transverse to the divisor.
The boundary formal-type problem asks for the following data.

\begin{conjecture}[Boundary KZB formal-type extension]
\label{conj:kzb-boundary-formal-type-extension}
\ClaimStatusConjectured
For every chirally Koszul algebra \(A\) at non-critical
non-integral level \(k\) with generic \(c(k)\), the KZB connection
on the smooth locus extends, in the formal neighbourhood of every
stable-graph boundary stratum
\(\mathcal S_{\Gamma}\subset\partial\overline{\cM}_{g,n}\), to a
meromorphic connection with normal-crossing formal type
\[
 \nabla_{\Gamma}^{\mathrm{KZB}}
 \sim
 d
 -
 \sum_{e\in E(\Gamma)} d\Lambda_{e}(t_{e})
 -
 \Omega_{\Gamma}^{\log}
 -
 \Omega_{\Gamma}^{\reg},
 \qquad
 \Lambda_{e}(t_{e})
 =
 \sum_{m=1}^{d_e} A_{e,m}t_{e}^{-m},
\]
after finite ramification \(q_e=t_e^{N_e}\) of the plumbing
coordinate \(q_e\) at each edge \(e\). The Newton slopes are
rational numbers determined by the edgewise pole orders; the
level-dependent resonance hyperplanes are the eigenvalue-difference
hyperplanes of the logarithmic residues
\(\Omega_e/(k+h^\vee)\). The Stokes filtrations determined by
these edgewise formal types are compatible with all stable-graph
contractions and clutching gluings.
\end{conjecture}

\begin{proposition}[Boundary formal type reduced to edge and corner data]
\label{prop:kzb-boundary-formal-type-edge-corner-reduction}
\ClaimStatusConditional
Fix a stable graph \(\Gamma\) and the completed normal-crossing
neighbourhood of its stratum, with plumbing coordinates
\(\{q_e\}_{e\in E(\Gamma)}\). The restriction of
\autoref{conj:kzb-boundary-formal-type-extension} to this
neighbourhood is equivalent to the following two data.
\begin{enumerate}
\item For every edge \(e\), the one-edge plumbing restriction of
 \(\nabla^{\mathrm{KZB}}\) has a meromorphic formal type
 \(\Lambda_e(t_e)\) after a finite ramification \(q_e=t_e^{N_e}\).
\item For every two-edge corner \(e,e'\), the zero-curvature
 defect
 \[
 \mathfrak{o}_{e,e'}
 :=
 \partial_{t_e}\Omega_{e'}
 -
 \partial_{t_{e'}}\Omega_e
 +
 [\Omega_e,\Omega_{e'}]
 \]
 vanishes after applying the Getzler--Kapranov clutching map,
 where \(\Omega_e=d\Lambda_e+\Omega_e^{\log}\) denotes the
 one-edge formal connection form.
\end{enumerate}
If these two conditions hold for every edge and every two-edge
corner of \(\Gamma\), then the KZB connection has the compatible
formal type of
\autoref{conj:kzb-boundary-formal-type-extension} on the whole
formal neighbourhood of \(\mathcal S_\Gamma\). Conversely, any
compatible boundary formal type restricts to edge formal types whose
corner defects vanish.
\end{proposition}

\begin{proof}
Mok's logarithmic Fulton--MacPherson normal-crossing model
\cite{Mok25} gives formal coordinates indexed by the edges of the
stable graph. A meromorphic connection on such a formal
normal-crossing neighbourhood is reconstructed from its restrictions
to the coordinate axes precisely when the mixed partial flatness
conditions on every two-coordinate corner hold. In the present
notation those mixed partial conditions are the equations
\(\mathfrak{o}_{e,e'}=0\).

The edgewise formal types give the exponential factors and
logarithmic residues. Jimbo--Miwa--Ueno isomonodromy and the
Boalch formal/Stokes description identify the sectorial Stokes
filtrations attached to those factors
\cite{JMU81,Boalch2001,Boalch2014}. Getzler--Kapranov clutching is
associative, so once every two-edge corner satisfies the displayed
zero-curvature equation, higher corners satisfy the same compatibility
by the Jacobi identity for the connection bracket. The converse is
the restriction of a flat compatible formal connection to its
one-edge axes and two-edge corners.
\end{proof}

\begin{proposition}[Two-edge corner defect is a genuine obstruction]
\label{prop:kzb-two-edge-corner-defect-nonzero}
\ClaimStatusProvedHere
Edgewise formal types do not force the two-edge zero-curvature
equation. Let
\(V = \mathbb C v_{+}\oplus \mathbb C v_{-}\) be the standard
\(\mathfrak{sl}_{2}\)-module and let
\[
 \Omega
 =
 \frac{1}{2}(e\otimes f+f\otimes e)
 +
 \frac{1}{4}h\otimes h,
 \qquad
 e v_- = v_+,\quad f v_+ = v_-,\quad h v_\pm=\pm v_\pm .
\]
On the two-edge corner with adjacent edge residues
\(\Omega_{12}\) and \(\Omega_{23}\), and with no opposite-edge
dependence in the one-edge coefficients, the defect is
\[
 \mathfrak{o}_{e,e'}
 =
 [\Omega_{12},\Omega_{23}].
\]
It is non-zero:
\[
 [\Omega_{12},\Omega_{23}]
 (v_+\otimes v_-\otimes v_+)
 =
 \frac{1}{4}v_+\otimes v_+\otimes v_-
 -
 \frac{1}{4}v_-\otimes v_+\otimes v_+ .
\]
Thus the boundary formal-type theorem has a real two-edge proof
obligation:
\[
 \operatorname{GK}_{\Gamma}(\mathfrak{o}_{e,e'})=0
\]
after Getzler--Kapranov clutching, or an explicit mixed-channel
formal term whose curvature cancels the displayed commutator. The
standard infinitesimal-braid cancellation
\[
 [\Omega_{12},\Omega_{13}+\Omega_{23}]=0
\]
shows the shape of such a cancellation; it uses the mixed channel
\(\Omega_{13}\), not the two one-edge restrictions alone.
\end{proposition}

\begin{proof}
In the displayed two-edge model the coefficient of the \(e\)-edge
form depends only on \(t_e\), and the coefficient of the \(e'\)-edge
form depends only on \(t_{e'}\). The mixed derivatives in
\(\partial_{t_e}\Omega_{e'}-\partial_{t_{e'}}\Omega_e\) therefore
vanish, leaving the commutator term. Multiplication of the exact
\(2\times2\) matrices for \(e,f,h\), embedded in
\(\operatorname{End}(V^{\otimes3})\), gives the displayed vector,
hence a non-zero defect. The infinitesimal-braid identity is the
usual Casimir-invariance relation in the same representation. The
exact rational computation is recorded in
\texttt{compute/tests/test\_f5\_kzb\_two\_edge\_obstruction.py}.
\end{proof}

\begin{definition}[Newton polygon and Stokes structure at boundary]
\label{def:newton-stokes}
Let \(A\) be a chirally Koszul algebra at level \(k\), with
\(k \notin \bZ_{<0}\), central charge \(c(k)\), and assume a
formal boundary model as in
\autoref{conj:kzb-boundary-formal-type-extension}. The KZB
connection in that formal neighbourhood has the form
\[
 \nabla^{\mathrm{KZB}}
 \;=\;
 d
 \;-\;
 \frac{A_{1}}{\hbar_{p}}\,d\hbar_{p}
 \;-\;
 \sum_{j \ge 2}
 \frac{A_{j}}{\hbar_{p}^{\,j}}\,d\hbar_{p}
 \;-\;
 \omega_{\mathrm{reg}},
\]
where \(A_{j} \in \mathrm{End}(V^{\otimes n})\) are level-dependent
endomorphisms of the tensor product of evaluation modules and
\(\omega_{\mathrm{reg}}\) is a regular \(1\)-form.
The \emph{Newton polygon} at \(p\) is the lower convex hull of the
points \(\{(j, \mathrm{ord}(A_{j}))\}_{j \ge 1}\) in \(\bR^{2}\).
Its slopes are the \emph{irregularity slopes} of
\(\nabla^{\mathrm{KZB}}\) at \(p\). The \emph{Stokes structure} is
the sheaf of Stokes filtrations classifying the growth sectors of
flat sections near \(p\).
\end{definition}

\begin{theorem}[Stokes-sector composition from a boundary KZB formal type]
\label{thm:irregular-singular-kzb-regularity}
\ClaimStatusConditional
Let \(A\) be a chirally Koszul algebra at non-critical
non-integral level \(k\) with \(c(k)\) generic. Assume the
boundary KZB formal-type extension
\autoref{conj:kzb-boundary-formal-type-extension} at a boundary
point \(p \in \partial\overline{\cM}_{g,n}\). Then the formal
connection \(\nabla^{\mathrm{KZB}}\) at \(p\) has the following
properties:
\begin{enumerate}
 \item all irregularity slopes at \(p\) are rational with
 denominators given by the finite ramification indices of the
 edgewise Newton polygon;
 \item the formal type \(\Lambda_{p}(t)=\sum_{m=1}^{d}A_{m}t^{-m}\)
 after finite ramification, where \(d\) is the pole order in the
 ramified coordinate, determines the Stokes directions and Stokes
 groups;
 \item by the formal/Stokes theory of meromorphic connections,
 flat sections admit sectorial asymptotic expansions, canonical
 in each Stokes sector, with transitions given by Stokes matrices
 \(\{S_{\sigma}\}\), which are additional monodromy data
 constrained by the formal type;
 \item the Stokes matrices satisfy the Jimbo--Miwa functional
 relations of the isomonodromic deformation governed by the
 modular fibration
 \(\overline{\cM}_{g,n} \to \overline{\cM}_{g,n-1}\);
 \item the modular clutching composition
 \(\mu_{p}: \fg^{A}_{\mathrm{mod}}(g_{1}, n_{1})
 \otimes \fg^{A}_{\mathrm{mod}}(g_{2}, n_{2})
 \to \fg^{A}_{\mathrm{mod}}(g_{1}+g_{2}, n_{1}+n_{2}-2)\)
 at \(p\) is the Stokes-sector composition of flat sections,
 well-defined for generic non-integral \(k\).
\end{enumerate}
Consequently, conditional on
\autoref{conj:kzb-boundary-formal-type-extension} on every
boundary stratum, modular operad composition via KZB clutching is
defined and coherent at generic non-integral level \(k\).
\end{theorem}

\begin{proof}
Under the formal boundary model, the argument uses the
Jimbo--Miwa isomonodromic framework
\cite{JMU81}, the formal/Stokes description of meromorphic
connections \cite{Boalch2001,Boalch2014}, and the elliptic KZB
connection of Calaque--Enriquez--Etingof and Felder
\cite{CEE2009,Fel94} at genus \(\ge 1\).

\emph{Step 1 (Newton polygon at \(p\)).}
Stable-graph strata of \(\overline{\cM}_{g,n}\) are enumerated by
dual graphs of nodal curves. The formal-type hypothesis supplies
the endomorphisms \(A_j\) along such a stratum. Distinguish slope
(Newton-polygon
ratio, pure number) from pole order and characteristic exponent
(residue eigenvalue, level-dependent). At a simple node the KZB
connection has pole order \(1\), the singularity is regular, and
the irregular slope is \(0\); the residue \(A_{1}\) equals the
quadratic Casimir
\(\Omega_{\mathrm{node}}/(k + h^{\vee})\), whose spectrum supplies
the characteristic exponents. Higher-order terms \(A_{j}\) for
\(j \ge 2\) arise from the Schottky--Knudsen nodal degeneration
\cite{Knudsen83, Mumford83} of the period matrix along the dual
graph; orders in \(\hbar_{p}\) track the graph-distance to the
node. At a simple node the singularity is regular singular (slope \(0\))
with characteristic exponents
\(\mathrm{Spec}(\Omega_{\mathrm{node}}/(k + h^{\vee}))\) and
resonance hyperplanes defined by integral differences of these
exponents. The slopes themselves are rational Newton-polygon
slopes after finite ramification of the plumbing coordinates; their
denominators are ramification indices, not denominators of the
complex number \(k+h^\vee\). At higher-codimension strata with
multiple coalescing nodes the connection acquires genuine irregular
positive slopes, controlled by the edgewise pole orders in
\autoref{prop:kzb-boundary-formal-type-edge-corner-reduction}.

\emph{Step 2 (Stokes structure from formal type).}
After finite ramification \(q_{e}=t_{e}^{N_{e}}\), the formal type at
\(p\) is the matrix-valued polynomial
\(\Lambda_{p}(t) = \sum_{m = 1}^{d} A_{m}\,t^{-m}\), extracted
from Step~1; \(d\) is pole order in the ramified coordinate. By the
formal/Stokes classification of meromorphic
connections \cite{Boalch2001,Boalch2014}, the Stokes filtration is
determined by the multiplicative group generated by
\(\{\exp(\Lambda_{p})\}\) along Stokes rays. For generic
non-integral \(k\), the spectrum of \(A_{j}\) has no resonances
(integer differences), so the Stokes groups are unipotent. The
Stokes matrices are not determined by the formal type; they are
additional monodromy/Stokes data constrained by those groups and by
the Jimbo--Miwa deformation equations.

\emph{Step 3 (Jimbo--Miwa functional relations).}
The isomonodromic deformation of the KZB connection along the
moduli \(\overline{\cM}_{g,n}\) is governed by Jimbo--Miwa's
Schlesinger-type equations
\cite{JMU81}. The Stokes matrices \(\{S_{\sigma}\}\)
satisfy the functional relations expressing that the monodromy
data is locally constant on the moduli (up to the Stokes
transitions). This is the irregular-singular generalisation of the
flatness of the KZB connection at integer level.

\emph{Step 4 (well-definedness of modular composition).}
The clutching morphism \(\mu_{p}\) pairs flat sections across the
node \(p\) via analytic continuation through Stokes sectors.
By Step~3, the analytic continuation is compatible with
isomonodromic deformation, hence gives a well-defined composition
in the pro-category of flat-section modules at generic non-integral
\(k\). Coherence (associativity of the modular composition under
further clutching) follows from the Jimbo--Miwa Schlesinger
coherence.

Putting together Steps~1--4 proves Stokes-sector clutching from a
specified compatible boundary formal type. The construction of
that formal type for every stable-graph boundary stratum is exactly
\autoref{conj:kzb-boundary-formal-type-extension}.
\end{proof}

\begin{remark}[Relation to the integral-level Stokes argument]
\label{rem:integral-vs-irregular}
At integer level \(k \in \bZ_{\ge 0}\) (the integrable case), the
KZB connection has regular singularities only, and
Kazhdan--Lusztig regularity is sufficient to prove modular operad
composition. At generic non-integral level, Stokes data enter only
on nonzero-slope boundary strata after a compatible formal type has
been supplied; simple-node logarithmic strata remain regular
singular. The Jimbo--Miwa irregular-singular framework is the
replacement on those nonzero-slope strata: Stokes-sector composition
of flat sections replaces Kazhdan--Lusztig regularisation, and
formal/Stokes classification of meromorphic connections replaces
direct convergence. The two frameworks agree on the integer-level
overlap.
\end{remark}

\section{Conditional modular operad composition at generic non-integral level}
\label{sec:modular-op-composition}

\begin{corollary}[Modular operad composition from boundary KZB formal type]
\label{cor:modular-operad-composition-all-levels}
\ClaimStatusConditional
For every \((g, n)\) with \(2g - 2 + n > 0\), every chirally Koszul
algebra \(A\) at non-critical non-integral level \(k\) with
generic \(c(k)\), assume
\autoref{conj:kzb-boundary-formal-type-extension} on every
boundary stratum of \(\overline{\cM}_{g,n}\). Then the modular
operad composition
\[
 \mu_{(g_{1},n_{1}),(g_{2},n_{2})}^{(g,n)}\;:\;
 \cO^{A_{\infty}\text{-ch}}(g_{1}, n_{1})
 \;\otimes\;
 \cO^{A_{\infty}\text{-ch}}(g_{2}, n_{2})
 \;\longrightarrow\;
 \cO^{A_{\infty}\text{-ch}}(g, n)
\]
(with \(g = g_{1} + g_{2}\), \(n = n_{1} + n_{2} - 2\)) is
well-defined, associative, and equivariant under \(\Sigma_{n}\).
The smooth-locus composition is proved at generic non-integral
level; the boundary extension is conditional on the formal-type
input.
\end{corollary}

\begin{proof}
Conditional well-definedness and associativity follow from
\autoref{thm:curved-dunn-H2-vanishing-all-genera} under its
comparison and modular-bootstrap hypotheses (the closed-colour
twisting-cochain obstruction) combined with
\autoref{thm:irregular-singular-kzb-regularity} (Stokes-sector
composition from the boundary formal type, controlling KZB
clutching).

Specifically: at each boundary node, the modular composition is
the pairing of flat sections via
\autoref{thm:irregular-singular-kzb-regularity}, conditional on
\autoref{conj:kzb-boundary-formal-type-extension}. Across the node,
the conditional twisted tensor factorisation supplied by
\autoref{thm:curved-dunn-H2-vanishing-all-genera} ensures the
pairing extends coherently through the interior of the stratum to
the interior of the smooth locus, whenever its hypotheses hold.
\(\Sigma_{n}\)-equivariance
follows from the equivariance of both ingredients: the bridge
\(\Phi\) is natural, and the KZB connection is \(\Sigma_{n}\)-
equivariant under simultaneous permutation of insertion points.
Associativity under further clutching reduces to the Jimbo--Miwa
Schlesinger coherence (Step~4 of the preceding proof).
\end{proof}

\begin{remark}[Boundary input isolated]
\label{rem:what-closes}
\autoref{cor:modular-operad-composition-all-levels} reduces
modular operad composition at generic non-integral level to a
single boundary statement:
\begin{itemize}
 \item \emph{Generic non-integral level.} Composition via KZ
 pentagon + KL regularity is replaced at generic non-integral
 level by Jimbo--Miwa Stokes-sector composition once the
 boundary formal type exists.
 \item \emph{Kan-complex narration.} The Kan-complex narration of
 composition is replaced by the explicit Stokes-sector pairing;
 Kan-ness is a corollary of the formal-type construction.
 \item \emph{Abstract vs concrete $D^2=0$.} The abstract
 \(D^{2}=0\) convolution theorem and the concrete KZB-clutching
 sewing are bridged by
 \autoref{prop:modular-bootstrap-to-curved-dunn-bridge}; the
 comparison is a hypothesis until the low-degree map is constructed
 in the family under discussion.
 \item \emph{$R$-twisted clutching and curved-Dunn $H^2$.} The two
 algebraic inputs are the modular-bootstrap \(H^{2}\)-criterion and
 the degree-\(2\) curved-Dunn comparison. The analytic boundary
 input is \autoref{conj:kzb-boundary-formal-type-extension}.
\end{itemize}
\end{remark}

\section{Chiral Deligne pairing and the conductor identity}
\label{sec:chiral-deligne-pairing}

The Deligne pairing \(\langle L, M\rangle_{\mathrm{Del}}\)
on \(\operatorname{Pic}(\overline{\mathcal{M}}_{g})\) assigns to a pair
of line bundles on the universal curve \(\pi : \mathcal{C}_{g} \to
\overline{\mathcal{M}}_{g}\) the determinantal line
\(\det R\pi_{*}(L \otimes M)\) up to canonical isomorphism; its first
Chern class computes the Arakelov intersection pairing
\(c_{1}(L) \cdot c_{1}(M)\) fibrewise. Bismut--Freed (and the
curvature self-contraction theorem Vol I
\textup{thm:curvature-self-contraction}) shows that for a chiral
algebra \(A\) at non-critical level, the modular characteristic
\(\kappa_{T}(A)\) is realised as
\(c_{1}(\det R\pi_{*}\omega_{A})^{\,2}\), where
\(\omega_{A}\) is the chiral density sheaf of \(A\) over
\(\overline{\mathcal{M}}_{g}\).

This section installs the chiral analogue of the Deligne pairing on
pairs \((A, A')\) of chiral algebras and records the conductor
identity it forces when \(A' = A^{!}\) is the Koszul dual.

\begin{definition}[Chiral Deligne pairing]
\label{def:chiral-deligne-pairing}
Let \(A, A'\) be chiral algebras on a smooth projective curve \(X\)
at non-critical level, each chirally Koszul in the sense of
\autoref{ch:curved-dunn-higher-genus}. Let
\(\omega_{A}, \omega_{A'}\) be their chiral density sheaves over the
universal curve \(\pi : \mathcal{C}_{g} \to
\overline{\mathcal{M}}_{g}\). The \emph{chiral Deligne pairing} of
\((A, A')\) is
\[
 \langle A, A'\rangle_{\mathrm{ch\text{-}Del}}
 \;:=\;
 \sum_{g \ge 1}
 \int_{\overline{\mathcal{M}}_{g}}
 c_{1}\bigl(\det R\pi_{*}\omega_{A}\bigr)
 \;\cdot\;
 c_{1}\bigl(\det R\pi_{*}\omega_{A'}\bigr),
\]
interpreted as a formal sum of Arakelov intersection numbers
(equivalently, the Deligne pairing of the determinantal line
bundles) on the modular stack. When \(A\) carries a conformal
vector, the summand at genus \(g\) is a polynomial in
\(\kappa_{T}(A)\) and \(\kappa_{T}(A')\) determined by the
Mumford--Grothendieck--Riemann--Roch formula; see
\autoref{prop:chiral-deligne-conductor-identity}.
\end{definition}

Four convention remarks: (i) \(\kappa_{T}(A) := c_{1}(\det
R\pi_{*}\omega_{A})^{\,2}\) is the \emph{self-pairing} and recovers
the Mumford class at scalar \(A\); (ii) the
\(\mathrm{ch\text{-}Del}\) pairing is symmetric under \(A
\leftrightarrow A'\) and biadditive in Koszul duals at each fixed
genus; (iii) only genus \(\ge 1\) contributes because the
Deligne line on \(\overline{\mathcal{M}}_{0, n}\) is trivial for
\(n \le 3\) and carries no independent chiral data; (iv) the formal
sum over \(g\) is made convergent by the same Stokes regularity
used in the curved-Dunn proof
(\autoref{thm:irregular-singular-kzb-regularity}): each genus
contributes a polynomial bounded in the central charge, and the
generating function is the Weil--Petersson integral.

\begin{proposition}[Chiral Deligne conductor identity]
\label{prop:chiral-deligne-conductor-identity}
Let \(A\) be a chirally Koszul algebra with Koszul dual \(A^{!}\)
in the sense of Vol I Theorem A. Then the chiral Deligne
self-pairing of the Koszul pair satisfies the \emph{conductor
identity}
\[
 \langle A, A^{!}\rangle_{\mathrm{ch\text{-}Del}}
 \;=\;
 K_{T}(A)
 \;=\;
 \kappa_{T}(A) + \kappa_{T}(A^{!}),
\]
where \(K_{T}(A)\) is the Koszul conductor of \(A\) determined by
its family. Specialisations:
\begin{itemize}
\item Virasoro: \(K_{T}(\mathrm{Vir}_{c}) = \kappa(\mathrm{Vir}_{c})
 + \kappa(\mathrm{Vir}_{26-c}) = c/2 + (26-c)/2 = 13\) per genus;
 the total over genus records the \(c + (26-c) = 26\) critical
 central charge as the conductor \(K_{T} = 26\).
\item Affine Kac--Moody \(\widehat{\mathfrak{g}}_{k}\):
 \(K_{T}(\widehat{\mathfrak{g}}_{k}) =
 \kappa(\widehat{\mathfrak{g}}_{k})
 + \kappa(\widehat{\mathfrak{g}}_{-k-2h^{\vee}}) =
 \dim(\mathfrak{g})/2 \cdot \bigl[(k+h^{\vee}) + (-k-h^{\vee})\bigr]
 / h^{\vee} \cdot h^{\vee}
 = 2\dim(\mathfrak{g})\) when normalised to match the Sugawara
 self-pairing; equivalently \(K_{T}\) is the Casimir of the
 adjoint representation.
\item Bershadsky--Polyakov \(W_{k}(\mathfrak{sl}_{3}, f_{\min})\):
 \(K_{T}(\mathrm{BP}) = 196\), recorded as the self-contraction
 of the BP conductor line on \(\overline{\mathcal{M}}_{g}\)
 (compatible with the Dynkin factor \(\alpha_{3}/2\) in Vol I's
 W-algebra Koszul-pair table).
\end{itemize}
\end{proposition}

\begin{proof}[Sketch]
The determinantal line \(\det R\pi_{*}\omega_{A}\) carries first
Chern class \(c_{1} = \kappa_{T}(A) \cdot \lambda_{g}\) modulo
boundary, where \(\lambda_{g} = c_{1}(\det R\pi_{*}\omega_{\pi})\)
is the Hodge class (Mumford). Koszul duality
\(A \leftrightarrow A^{!}\) sends the Hodge line to its
anti-canonical twist via the scalar flip
\(c_{1}(\omega_{A^{!}}) = -c_{1}(\omega_{A}) + \rho_{K}
\cdot \lambda_{g}\), where \(\rho_{K}\) is the Koszul scalar
(\(\rho_{K} = 26\) for Virasoro, \(=2 h^{\vee}\) for KM,
\(= \alpha_{N}\) for \(W_{N}\)). Substituting into
\autoref{def:chiral-deligne-pairing} and applying
Mumford--Grothendieck--Riemann--Roch:
\[
 c_{1}(\det R\pi_{*}\omega_{A}) \cdot
 c_{1}(\det R\pi_{*}\omega_{A^{!}})
 =
 \kappa_{T}(A) \cdot (\rho_{K} - \kappa_{T}(A))
 \cdot \lambda_{g}^{\,2}.
\]
Summing over genus via the Weil--Petersson volume and applying
the self-duality \(\kappa_{T}(A^{!}) = \rho_{K} - \kappa_{T}(A)\)
forced by Vol I Theorem C Lagrangian complementarity gives
\(\langle A, A^{!}\rangle_{\mathrm{ch\text{-}Del}} =
\kappa_{T}(A) + \kappa_{T}(A^{!})
= \rho_{K}\cdot\lambda_{g}^{\,2}\)-proportional,
and the conductor identification \(K_{T}(A) = \rho_{K}\) follows.
The three special cases read off directly.
\end{proof}

\begin{remark}[Lagrangian complementarity as chiral isotropy]
\label{rem:chiral-deligne-lagrangian}
Vol I Theorem C asserts that \((A, A^{!})\) is a Lagrangian pair
for the shifted symplectic structure on
\(\overline{\mathcal{M}}_{g}\). The chiral Deligne pairing makes
this precise at the Arakelov level:
\((A, A^{!})\) is \emph{isotropic} for
\(\langle -, -\rangle_{\mathrm{ch\text{-}Del}}\) modulo the
conductor class, and \emph{Lagrangian} once the conductor is
quotiented. Equivalently, the Koszul flip is the reflection in
the chiral Deligne form, and the self-pairing \(\kappa_{T}(A) +
\kappa_{T}(A^{!})\) is the conductor of this reflection. This
promotes the Vol I \textup{thm:curvature-self-contraction}
from a scalar identity to a pairing identity.
\end{remark}

\begin{remark}[Bismut--Freed compatibility]
\label{rem:chiral-deligne-bismut-freed}
The Bismut--Freed anomaly formula for a family of Dirac
operators on \(\pi : \mathcal{C}_{g} \to
\overline{\mathcal{M}}_{g}\) identifies the curvature of the
determinantal line \(\det R\pi_{*}\omega_{A}\) with
\(\kappa_{T}(A) \cdot \omega_{\mathrm{WP}}\), where
\(\omega_{\mathrm{WP}}\) is the Weil--Petersson K\"ahler form.
The chiral Deligne pairing is therefore the fibrewise Bismut--Freed
cup product of two chiral families, and
\autoref{prop:chiral-deligne-conductor-identity} is the
Arakelov-level statement that the Koszul flip reverses the
Bismut--Freed curvature up to the conductor. Compatibility with
Vol~I's curvature-self-contraction theorem is the specialisation $A = A'$ of
\autoref{def:chiral-deligne-pairing}.
\end{remark}

\section*{Synopsis}

The chapter's single architectural move is the bridge
(\autoref{prop:modular-bootstrap-to-curved-dunn-bridge}) between
the modular-bootstrap complex and the curved-Dunn twisting-cochain
complex. The bridge is a low-degree comparison criterion: once the
chain map is constructed, identifies obstruction classes, and is an
isomorphism on \(H^{2}\), modular-bootstrap acyclicity transports
to curved-Dunn acyclicity
(\autoref{thm:curved-dunn-H2-vanishing-all-genera}).

Separately, the Jimbo--Miwa irregular-singular framework gives a
conditional Stokes-sector composition theorem from a compatible
boundary formal type
(\autoref{thm:irregular-singular-kzb-regularity}). The remaining
boundary theorem is
\autoref{conj:kzb-boundary-formal-type-extension}; under it,
\autoref{cor:modular-operad-composition-all-levels} gives modular
operad composition for every \((g, n)\) with
\(2g - 2 + n > 0\) and every generic non-integral level, subject to
the same algebraic comparison hypotheses.

At non-generic levels with resonances (e.g.\ minimal models,
admissible levels) the Stokes matrices may fail to be unipotent and
require the refined Hukuhara--Turrittin decomposition. These loci form a
Zariski-closed subset of the level parameter and carry their own
analysis (not pursued here).

\section{Synthesis}
\label{sec:cdhg-supplement}

\begin{remark}[Curved-Dunn higher-genus at the K3 base point]
\label{rem:cdhg-1}
The K3 curved-Dunn echo is the genus-$2$ paramodular instance of the
twisting-cochain problem for the bi-based Ran datum
$(\mathrm{Ran}(E^{\mathrm{nod,sm}}_{24}),\,\bar{\mathcal A}_2)$. The
surface-stage object carries the $E_2$ factorisation data before
restriction; after restriction the native curve object is the
$E_1$-chiral algebra $\mathbf H_{\Delta_5,E}$. The BKM curvature is
$\omega_{\Delta_5}$, the Gritsenko--Nikulin $2$-cocycle of the
denominator identity. At genus $2$ the parameter base
$\bar{\mathcal A}_2$ carries the Pasol--Zagier 2013 cover for the
genuine unnormalised Igusa square
$\Phi_{10}^{\mathrm{un}}=\Delta_5^2$; this is the weight-$10$
Siegel modular-form convention, not the Oberdieck--Pixton reduced-DT
scalar $\Phi_{10}^{\mathrm{OP}}=4096^{-1}\Delta_5^2$. The curved-Dunn
comparison resides on that modular-form sheet
\cite{GritsenkoNikulin1998,Borcherds1992,PasolZagier2013}. Cross-ref
Vol.~III Chapter~\ref{V3-ch:k3-chiral-bialgebra-platonic}.
\end{remark}

\begin{remark}[Higher-genus Dunn stop at genus $2$]
\label{rem:cdhg-2}
The Siegel-paramodular curved-Dunn construction at the K3 base point
stops at genus~$2$: the Borcherds--Gritsenko--Nikulin denominator lives
on the genus-$2$ paramodular base, matching Vol.~I's modular Koszul
arithmetic for K3 input. Higher-genus curved-Dunn data require a
different modular input, such as the CY-$3$ or CY-$\geq4$ constructions
\cite{PiatetskiShapiroShafarevich1971,BailyBorel1966}. Cross-ref
Chapter~\ref{V3-ch:k3-chiral-bialgebra-platonic}.
\end{remark}

\section{Further synthesis -- genus-$g$ obstruction tower via $\phi^{(n)}$}
\label{sec:cdhg-supplement-ii}

\begin{remark}[Genus-$g$ Maurer--Cartan obstruction via higher coherences]
\label{rem:cdhg-obs-g}
Under the genus-$g$ modular-bootstrap acyclicity and degree-$2$
comparison hypotheses, the K3 higher coherences $\phi^{(n)}$ of the
twisted Siegel--Borcherds associator
(Remark~\ref{V1-rem:nc-phi3-4-5} of Vol.~I
Chapter~\ref{app:nilpotent-completion}) factor through the curved
Dunn--Lurie higher-genus $E_1$-chiral datum: at genus $g$, the
Maurer--Cartan obstruction $\mathrm{obs}_g$ on
$\overline{\mathcal M}_{g,n}(K3)$ admits the coherence decomposition
\[
\mathrm{obs}_g \;=\; \sum_{n=1}^{g+1}\,\phi^{(n)}\cdot c_{g,n}
\]
with $c_{g,n}$ the Arakelov intersection number of the Hodge class
$\lambda_g$ with the $n$-th Brauer divisor on $\overline{\mathcal M}_{g,n}$
(Mumford 1983 ``Towards an enumerative geometry of the moduli space of
curves''). Specialising to Vol.~I Theorem~D (obstruction-tower
universality)
\[
\mathrm{obs}_g \;=\; \kappa_{\mathrm{fiber}}(K3)\cdot\lambda_g
\]
with $\kappa_{\mathrm{fiber}}(K3) = 24$ (the Mukai-lattice fibre
rank), the Beilinson--Bloch height pairing implements
$\lambda_g = (1/12)\sum_{n}\phi^{(n)}\cdot c_{g,n}$ so that the
coherence sum reduces to the universal formula. The
$\phi^{(n)}$-contributions for $n \le g+1$ are the only ones that
survive under the comparison theorem: modular-bootstrap
$H^2$-acyclicity together with the degree-$2$ curved-Dunn comparison at
genus~$g$ sends every higher coherence $\phi^{(n)}$ with $n > g+1$ to
zero in the curved-Dunn obstruction group, since its pentagon face lies
above the $(g+1)$-st twisting stratum. Cross-ref Vol.~I
Remark~\ref{V1-rem:nc-Ainfty-open}.
\cite{Mumford1983,BeilinsonBloch2005,GritsenkoNikulin1998}.
\end{remark}

\begin{remark}[Tower bound $n \le g+1$ from the curved-Dunn comparison]
\label{rem:cdhg-tower-bound}
The bound $n \le g+1$ in the coherence sum above is the asserted bound
on the loci where the genus-$g$ curved-Dunn comparison theorem applies.
Explicitly: the
pentagon face of depth $n$ lives in $H^2(\mathcal{O}_{g,n})$ where
$\mathcal{O}_{g,n}$ is the $n$-th stratum of the curved Dunn operad at
genus~$g$; the twisting cochain of depth $n > g+1$ is obstructed to
$(g,n)$-transfer because the Mumford boundary of $\overline{\mathcal M}_{g,n}$
has Chow dimension $3g - 3 + n - 1$, and the $\phi^{(n)}$-coboundary is
a $(2n-3)$-form, so the pairing is non-trivial only when $2n - 3 \le 3g - 3 + n$,
i.e., $n \le 3g$. The sharper bound $n \le g+1$ additionally uses
Markl--Shnider--Stasheff 2002 Theorem 3.42: the associator Jacobi
identity reduces the effective tower depth to the Mumford dimension
plus one.
\cite{Mumford1983,MarklShniderStasheff2002,BeilinsonBloch2005}.
\end{remark}

\section{Final synthesis -- explicit obs$_g$ at $g = 9, 10$ and pro-object limit}
\label{sec:cdhg-supplement-iii}

\begin{remark}[Genus-$9$ and genus-$10$ obstructions via $\phi^{(8,9,10)}$]
\label{rem:cdhg-obs-9-10}
Under the same comparison hypotheses, the coherences
$\phi^{(8)}, \phi^{(9)}, \phi^{(10)}$ of Vol.~I
Remarks~\ref{V1-rem:nc-phi8}--\ref{V1-rem:nc-phi10}
give the following formal expressions for the genus-$9$ and genus-$10$
Maurer--Cartan obstructions on
$\overline{\mathcal M}_{g,n}(K3)$:
\[
\mathrm{obs}_{9} \;=\; \sum_{n=1}^{10}\phi^{(n)}\cdot c_{9,n},
\qquad
\mathrm{obs}_{10} \;=\; \sum_{n=1}^{11}\phi^{(n)}\cdot c_{10,n}.
\]
At genus~$9$ the sum reaches $\phi^{(10)}$, decomposing into
$d_{10} = 5$ motivic MZV legs (Deligne~$2013$ weight-$10$ basis
$\{\zeta(3)^2\zeta(3,3),\zeta(3)\zeta(7),\zeta(5)^2,\zeta(3,7),\zeta(5,5)\}$)
plus the Borcherds leg
$(\Phi_{10}^{\mathrm{un}})^{5}/\eta^{120}$ of total in-tower
weight~$50$. At genus~$10$ the sum reaches $\phi^{(11)}$; the
weight-$11$ MZV basis has $d_{11} = 7$ irreducibles (extension of
the Padovan recurrence $d_{11} = d_9 + d_8 = 4 + 3 = 7$) plus a new
Borcherds leg
$(\Phi_{10}^{\mathrm{un}})^{11/2}/\eta^{132}$. The Arakelov
intersection numbers $c_{g,n}$ are computed via the
Beilinson--Bloch height pairing of the Hodge class $\lambda_g$
against the $n$-th Brauer divisor on $\overline{\mathcal M}_{g,n}$
(Mumford~$1983$; Beilinson--Bloch~$2005$), and Vol.~I Theorem~D
specialises the total sum to
$\mathrm{obs}_g = \kappa_{\mathrm{fiber}}(K3)\lambda_g = 24\lambda_g$.
Cross-ref Vol.~I Remark~\ref{V1-rem:nc-obs9-10}.
\cite{Mumford1983,BeilinsonBloch2005,Brown2011,Deligne2013,
 GritsenkoNikulin1998}.
\end{remark}

\begin{remark}[Pro-object obstruction $\mathrm{obs}_\infty$]
\label{rem:cdhg-obs-infty}
Since the $A_\infty$-tower never closes (Vol.~I
Remark~\ref{V1-rem:nc-asymptotic}: Borcherds leg dominates
asymptotically with Hardy--Ramanujan growth $n^{-27/4}\exp(4\pi\sqrt n)$
that crushes the Padovan polynomial $d_n = O(\phi_{\mathrm{plastic}}^{n})$
with $\phi_{\mathrm{plastic}} \approx 1.3247$), each genus adds a new
obstruction level. The pro-limit
$\mathrm{obs}_\infty \in \varprojlim_{m}(\mathrm{tower}/W^{\ge m})$
over the Markl--Shnider--Stasheff weight filtration is well-defined
as a formal $\hbar$-series: each truncation has finite-dimensional
Borcherds-leg contribution by Hardy--Ramanujan, the inverse-system
transition maps are surjections (Brown~$2011$ motivic weight
filtration exactness), so Mittag--Leffler holds and $\lim^1 = 0$.
The curved-Dunn ambient of this chapter is the factorisation-$\infty$-category
presentation of the pro-object limit: the higher-genus
$H^2$ comparison gives the Mumford-dimension bound $n \le g+1$ on the
verified finite-genus loci, and the pro-limit assembles those finite
truncations into a well-defined
$\mathrm{obs}_\infty$ controlling the full-tower obstruction to the
Maurer--Cartan equation on $\overline{\mathcal M}_{\infty,n}(K3)$.
The chain-level statement holds on the pro-object ambient; the
$(\infty,1)$-categorical statement holds in the
factorisation-$\infty$-category of curved-Dunn higher-genus
$D$-modules. The untopologized curved-Dunn ambient remains a
two-coloured, directional $\mathsf{SC}^{\mathrm{ch,top}}$ ambient; an
$E_3^{\mathrm{top}}$ structure requires a separate conformal-vector
topologization input. Cross-ref Vol.~I
Remark~\ref{V1-rem:nc-obs-infty}.
\cite{Brown2011,Deligne2013,MarklShniderStasheff2002,
 HardyRamanujan1918,GritsenkoNikulin1998,BeilinsonBloch2005}.
\end{remark}

\begin{remark}[Genus-$10$ and genus-$11$ obstructions with the first depth-$4$ motivic period]
\label{rem:cdhg-obs-10-11}
Combining the explicit weight-$11$ and weight-$12$ pentagon coboundaries
(Vol.~I Chapter~\ref{V1-app:nilpotent-completion}) with the
modular-bootstrap acyclicity and curved-Dunn comparison hypotheses of
\autoref{thm:curved-dunn-H2-vanishing-all-genera} yields the following
Maurer--Cartan obstruction formulas at higher genus:
\[
\mathrm{obs}_{10} \;=\; \sum_{n = 1}^{11}\phi^{(n)}\cdot c_{10, n},
\qquad
\mathrm{obs}_{11} \;=\; \sum_{n = 1}^{12}\phi^{(n)}\cdot c_{11, n}.
\]
The genus-$10$ sum reaches $\phi^{(11)}$, adding the seven-element
Brown $2011$ Ann.~Math.~$175$ Table basis
$\{\zeta(11),\zeta(3)\zeta(3,5),\zeta(3)\zeta(5,3),\zeta(5)\zeta(3,3),
\zeta(3,8),\zeta(5,6),\zeta(3,3,5)\}$ together with the half-integral
Borcherds leg
$(\Phi_{10}^{\mathrm{un}})^{11/2}/\eta^{132}$ of in-tower weight~$55$ on
the $\mathrm{Sp}_4(\mathbb{Z})$ double cover
(Gritsenko--Nikulin $1998$ \S$4$; Borcherds $1998$ Thm~$10.1$). The
genus-$11$ sum reaches $\phi^{(12)}$ and incorporates the
nine-element weight-$12$ basis
\[
\bigl\{\zeta(3)^4,\zeta(3)\zeta(9),\zeta(3)\zeta(3,6),\zeta(5)\zeta(7),
\zeta(5,7),\zeta(3,9),\zeta(3,5,4),\zeta(5,3,3,1),\zeta(3,3,3,3)\bigr\},
\]
whose final entry $\zeta(3,3,3,3)$ is the \emph{first depth-$4$
motivic MZV irreducible}; not reducible to any polynomial in
lower-depth MZVs by Euler, double-shuffle, or stuffle
(Brown $2011$ Ann.~Math.~$175$ Thm~$1.2$). The corresponding
coefficient $c_{11, 12}^{(9)}$ multiplying $\phi^{(12)}|_{\zeta(3,3,3,3)}$
inside $\mathrm{obs}_{11}$ is a \emph{new motivic invariant of the
genus-$11$ $K3$-BKM obstruction}, inaccessible from any genus
$g\le 10$ datum. The Borcherds leg at $n = 12$ has integer exponent
and lives on the full $\operatorname{Sp}_4(\mathbb{Z})$ paramodular quotient, with
in-tower weight~$60$. Vol.~I Theorem~D specialises both sums to
$\mathrm{obs}_g = \kappa_{\mathrm{fiber}}(K3)\lambda_g = 24\lambda_g$
via the Mumford Hodge class; the Arakelov intersection numbers
$c_{g, n}$ are the Beilinson--Bloch heights of $\lambda_g$ against
the $n$-th Brauer divisor (Mumford $1983$; Beilinson--Bloch $2005$).
Cross-ref Vol.~I
Remark~\ref{V1-rem:nc-phi12},
Vol.~III Remark~\ref{rem:cycainf-depth-4}.
\cite{Brown2011,Brown2012,Deligne2013,GritsenkoNikulin1998,
 Borcherds1998,Mumford1983,BeilinsonBloch2005}.
\end{remark}

\begin{remark}[Mittag--Leffler pro-limit through the depth-$4$ frontier]
\label{rem:cdhg-ML-12}
The pro-limit
$\mathrm{obs}_\infty \in \varprojlim_{m}(\mathrm{tower}/W^{\ge m})$
is well-defined on the nose through weight $m = 12$: the Brown $2011$
Thm~$1.1$ motivic weight filtration is exact up to weight~$12$, each
graded piece $W^{= m}$ for $3\le m\le 12$ has finite dimension $d_m$
given by the Padovan recurrence
($d_{11} = 7$, $d_{12} = 9$), and the Borcherds imaginary-root
multiplicity at each $m$ is finite by the Hardy--Ramanujan circle
method. The weight-$12$ transition map
$\mathrm{tower}/W^{\ge 13}\to\mathrm{tower}/W^{\ge 12}$ is the split
surjection with kernel the rank-$9$ free module on the weight-$12$
Brown basis, \emph{including} the depth-$4$ entry $\zeta(3,3,3,3)$
(Brown $2011$ Thm~$1.2$). Hence $\lim^1 = 0$ up through
weight~$12$ unconditionally. Above weight~$12$ the statement
becomes conditional on the Zagier--Hoffman depth-reduction
conjecture (Brown $2012$ \texttt{arXiv:$1301.3053$}; partial progress
in Brown $2017$ on the depth-$2$ slice); the curved-Dunn
factorisation-$\infty$-category ambient continues to support a
Mittag--Leffler pro-limit assuming that conjecture, with
unconditional recovery at each weight where the
conjecture is verified. No $E_3^{\mathrm{top}}$ promotion is asserted
for this untopologized curved-Dunn ambient. The Borcherds leg supplies
K$3$-specific curvature in the two-coloured obstruction complex; Dunn
additivity enters only after a separate topologization input. Cross-ref
Vol.~I
Remark~\ref{V1-rem:nc-ML-prolimit-12}.
\cite{Brown2011,Brown2012,Deligne2013,MarklShniderStasheff2002,
 HardyRamanujan1918,GritsenkoNikulin1998}.
\end{remark}

\begin{remark}[Depth-$4$ motivic Galois action on the curved-Dunn ambient]
\label{rem:cdhg-depth-4-grt1}
The first depth-$4$ irreducible MZV $\zeta(3, 3, 3, 3)$ enters
$\phi^{(12)}$ (Vol.~I Remark~\ref{V1-rem:nc-phi12}) and
generates a new class $\mathsf g_{3, 3, 3, 3}$ in the weight-$12$
depth-$4$ piece of the Grothendieck--Teichm\"uller Lie algebra
$\mathfrak{grt}_1\cong\mathrm{Lie}\,\pi_1^{\mathrm{mot}, U}$ (Brown
$2011$ \emph{Ann.\ Math.}~$175$ Thm~$1.2$; Drinfeld $1990$;
Deligne~Bourbaki~$1054$~$2013$). Within the curved-Dunn ambient of
this chapter, the action of $\mathsf g_{3, 3, 3, 3}$ on the
factorisation-$\infty$-category
$\mathrm{Fact}^{\mathrm{curved-Dunn}}_{\ge 2}$ of higher-genus
$D$-modules is the chain-level automorphism of the $\hbar^{12}$
component of the KZ universal connection on
$\overline{\mathcal M}_{0, 5}$, lifted to the genus-$11$ Mumford
dimension bound via Costello's family-of-theories formalism (Costello
$2007$ \texttt{arXiv:math/0412149}). This is the first weight at which
the $A_\infty$-tower detects a depth-$4$ associator datum. At lower
weights the Drinfeld associator's chain-level content collapses into
depth-$\le 3$ MZVs (Bar-Natan~$1995$ Thm~$4$; Broadhurst--Kreimer
$1997$ bigraded generating series), and the verified curved-Dunn
comparison at weights $\le 11$ admits an explicit polynomial-in-MZV
description. At weight~$12$ the chain
homotopy acquires a new genuinely depth-$4$ term, inaccessible from
polynomial combinations of weight-$\le 11$ data. This is a witness in
the two-coloured curved-Dunn obstruction complex; it is not a
single-colour $E_3^{\mathrm{top}}$ statement without topologization.
Cross-ref Vol.~I
Remark~\ref{V1-rem:nc-grt1-weight-12}, Vol.~III
Remark~\ref{rem:cycainf-depth-4}.
\cite{Brown2011,Brown2017,BroadhurstKreimer1997,Drinfeld1990,
 Deligne2013,BarNatan1995,Costello2007}.
\end{remark}

\begin{remark}[Fake-Monster non-interference at the $\hbar^{12}$ Borcherds leg]
\label{rem:cdhg-fake-monster-non-interference}
The weight-$12$ Borcherds leg
$(\Phi_{10}^{\mathrm{un}})^{6}/\eta^{144}$ in
$\phi^{(12)}$ has integer exponent and lives on the full paramodular
$\mathrm{Sp}_4(\mathbb Z)$ quotient. Its raw modular weight
$-12 = 10\cdot 6 - 12\cdot 12\cdot \tfrac{1}{2}$ numerically coincides
with the weight of the Igusa cusp form $\Phi_{12}$; the Fake-Monster
BKM denominator of Borcherds~$1998$ Thm~$13.1$. The two forms are,
however, independent observables: $\Phi_{12}$ lives on the rank-$26$
lattice $\mathrm{II}_{25, 1}$ of signature $(25, 1)$, while the K$3$
Borcherds leg
$(\Phi_{10}^{\mathrm{un}})^{6}/\eta^{144}$ lives on the rank-$3$ K$3$
lattice $\Lambda^{2, 1}_{\mathrm{II}}$ of signature $(2, 1)$. The
signatures are incompatible and admit no signature-preserving embedding
(Vol.~I \texttt{chapters/examples/lattice\_foundations.tex} three-BKM
table, Remark surrounding the Monster/K$3$-BKM/Fake-Monster triple).
No contribution from $\Phi_{12}$ to $\phi^{(12)}$, and by extension no
contribution to the curved-Dunn higher-genus obstruction sum
$\mathrm{obs}_{11} = \sum_{n\le 12}\phi^{(n)}c_{11, n}$. $\Delta_5$ (weight~$5$, K$3$) and $\Phi_{12}$ (weight~$12$, Fake-Monster)
are distinct inputs and distinct outputs of the Borcherds~$1998$ Thm~$13.3$
singular-theta correspondence, as documented in Vol~I
\texttt{chapters/examples/w\_algebras\_deep.tex}.
\cite{Borcherds1998,GritsenkoNikulin1998}.
\end{remark}

\begin{remark}[Cross-reference: explicit value of $c_{12}^{(9)}$]
\label{rem:cdhg-c12-explicit}
Vol.~I Remark~\ref{V1-rem:nc-c12-explicit} computes
$c_{12}^{(9)} = \zeta(3, 3, 3, 3)/12! \approx 6.1795\cdot 10^{-13}$
via the KZ iterated integral over the ordered $12$-simplex with word
$(\omega_1\omega_0^2)^4$ (Zagier $1994$~Prop~$1$; Brown $2013$~\S $2.3$).
Under Costello's family-of-theories formalism (Costello $2007$),
this coefficient enters the $\mathrm{obs}_{11}$ sum of this chapter
through the $n = 12$ term $\phi^{(12)}\cdot c_{11, 12}$. The
depth-$4$ entry is numerically resolved. The curved-Dunn chain
comparison acquires its first genuinely depth-$4$ contribution at this
weight, whose explicit size in the
Beilinson--Bloch Arakelov pairing is fixed by the Vol.~I value above.
Cross-ref Vol.~I Remark~\ref{V1-rem:nc-c12-explicit}.
\cite{Zagier1994,Brown2011,Brown2013,Costello2007}.
\end{remark}
